\chapter{貪心}
\label{ch:greedy}

% 本章的圖共用的樣式（前綴 gr-）
\tikzset{
    gr-leaf/.style={rectangle, draw=black!65, fill=coreFill, align=center,
        font=\scriptsize, inner sep=1.6pt, minimum width=0.55cm},
    gr-inner/.style={circle, draw=black!65, fill=white, minimum size=0.6cm,
        inner sep=0pt, font=\scriptsize},
    gr-new/.style={draw=memStroke, very thick, fill=memFill},
    gr-dummy/.style={rectangle, draw=black!45, dashed, fill=white, align=center,
        font=\scriptsize, inner sep=1.6pt, minimum width=0.55cm},
    gr-edge/.style={draw=black!60},
    gr-bit/.style={font=\tiny, text=black!60, inner sep=1pt},
    gr-plab/.style={font=\small},
    gr-mark/.style={draw=black!45, dashed},
}

系上的小會議室明天收到十二筆借用申請，每一筆寫著開始和結束的時刻，有些時段互相重
疊。 會議室一次只能給一筆申請用，系辦想讓借得到的申請越多越好。 有人說先排最早開始
的，有人說先排用得最短的，也有人說先排最早結束的，排進去之後，和它衝突的就不要了。
這三種做法哪一種一定能排到最多筆？ 往下讀之前，先想一想，再試著各舉一個反例。

上面三種做法都是\textbf{貪心}（greedy）：每一步只看眼前，挑一個「現在看起來最好」
的選擇，做了就不再反悔，然後去處理剩下的問題。 這樣的演算法通常很短、很快，難的是
證明它對。 很多看起來合理的貪心是錯的，所以這一章每一個貪心演算法都會附上證明，最
常用的工具是\textbf{交換論證}（exchange argument）：拿任何一個最佳解，一步一步把它
換成貪心的選擇，證明每一步都不會變差。

第~\ref{sec:gr-activity}~節用活動選擇示範交換論證，第~\ref{sec:gr-knapsack}~節比較分數背包和
0/1 背包，第~\ref{sec:gr-huffman}~節講霍夫曼編碼和外部排序的最佳合併樹，
第~\ref{sec:gr-sched}~節是兩個排程問題（第二個用互斥集合加速），第~\ref{sec:gr-fail}~節
看貪心什麼時候不對，並介紹擬陣。

這一章用到第~8~章的最小堆積、第~13~章的互斥集合、第~14~章的排序和第~15~章的外部排
序。 第~18~章的動態規劃和這一章常常互相對照：貪心解不了的問題，很多要改用動態規劃。
C 程式用到《C 程式設計基礎》第~14~章的結構、第~15~章的遞迴、第~18~章的
\texttt{qsort} 和第~19~章的 \texttt{malloc}，用到時先在框裡複習。

\section{活動選擇}
\label{sec:gr-activity}

\subsection{問題}

有 $n$ 筆申請，第 $i$ 筆要用的時段是半開區間 $[s_i, f_i)$，$s_i < f_i$：從時刻 $s_i$
開始，到時刻 $f_i$ 結束。 兩筆申請的時段沒有交集，就說它們\textbf{相容}
（compatible）。 因為是半開區間，一筆在 $f_i$ 結束，另一筆在同一個時刻 $f_i$ 開始，
不算衝突。 \textbf{活動選擇問題}（activity-selection problem）要挑出最多筆兩兩相容的
申請。

章首的十二筆申請如下（時刻用 $24$ 小時制的整點）：

\begin{center}
\begin{tabular}{c*{12}{c}}
\toprule
編號 & 1 & 2 & 3 & 4 & 5 & 6 & 7 & 8 & 9 & 10 & 11 & 12 \\
\midrule
$s_i$ & 13 & 9 & 17 & 8 & 12 & 10 & 15 & 8 & 11 & 16 & 14 & 19 \\
$f_i$ & 16 & 11 & 20 & 13 & 13 & 12 & 18 & 10 & 14 & 19 & 15 & 21 \\
\bottomrule
\end{tabular}
\end{center}

\subsection{三種貪心，兩種會錯}

章首的三種做法都是「依某個順序看申請，和已選的相容就選」，差別只在順序。 先看兩
個錯的。

\textbf{最早開始的先排。} 申請 $[8, 18)$、$[9, 10)$、$[11, 12)$：最早開始的是
$[8, 18)$，選了它，另外兩筆都和它衝突，只排到 $1$ 筆；不選它可以排 $2$ 筆。 一筆很早
開始、拖很久的申請會把後面整段時間都佔掉。

\textbf{用得最短的先排。} 申請 $[9, 12)$、$[11, 13)$、$[12, 15)$：最短的是
$[11, 13)$，兩小時，可是它同時和另外兩筆衝突。 選它只排到 $1$ 筆；選另外兩筆可以排
$2$ 筆。

\textbf{最早結束的先排。} 這一個是對的。 直覺是：先結束的申請，把會議室還回來得最
早，留給後面的時間最多。 下一小節把這個直覺變成證明。

「和別人衝突最少的先排」也會錯，只是反例要十來筆申請才造得出來。

\subsection{演算法}

先把申請依結束時刻由小到大排好，再從頭看一次：和最後選進來的那一筆相容，就選。
因為選進來的申請結束時刻是遞增的，「和已選的全部相容」只要檢查「開始時刻不早於最
後一筆的結束時刻」。

\begin{pseudo}[label=alg:gr-activity]{活動選擇（申請已依 $f$ 由小到大排好，$n \ge 1$）}
\Function{Activity-Select}{$s, f, n$}
    \State $A \gets \{1\}$
    \State $\mathit{last} \gets f[1]$ \Comment{最後選進來的那筆的結束時刻}
    \For{$i \gets 2$ \To $n$}
        \If{$s[i] \ge \mathit{last}$}
            \State $A \gets A \cup \{i\}$
            \State $\mathit{last} \gets f[i]$
        \EndIf
    \EndFor
    \State \Return $A$
\EndFunction
\end{pseudo}

排序 $\Theta(n \lg n)$，第 4--7 行的迴圈 $\Theta(n)$，總共 $\Theta(n \lg n)$。 輸入已經
依結束時刻排好的話只要 $\Theta(n)$。 結束時刻相同的申請，本書依編號由小到大排；不管
怎麼排，選到的筆數都一樣（觀念題第~\ref{q:gr-act-trace}~題可以自己驗證）。

\subsection{手算追蹤}

章首的十二筆依結束時刻排好之後逐筆檢查，$\mathit{last}$ 是檢查之前的值：

\begin{center}
\begin{tabular}{cccccl}
\toprule
編號 & $[s_i, f_i)$ & $\mathit{last}$ & $s_i \ge \mathit{last}$？ & 結果 & 選進來之後 $\mathit{last}$ \\
\midrule
8 & $[8, 10)$ & -- & -- & 選（第一筆） & 10 \\
2 & $[9, 11)$ & 10 & $9 < 10$ & 跳過 & \\
6 & $[10, 12)$ & 10 & $10 \ge 10$ & 選 & 12 \\
4 & $[8, 13)$ & 12 & $8 < 12$ & 跳過 & \\
5 & $[12, 13)$ & 12 & $12 \ge 12$ & 選 & 13 \\
9 & $[11, 14)$ & 13 & $11 < 13$ & 跳過 & \\
11 & $[14, 15)$ & 13 & $14 \ge 13$ & 選 & 15 \\
1 & $[13, 16)$ & 15 & $13 < 15$ & 跳過 & \\
7 & $[15, 18)$ & 15 & $15 \ge 15$ & 選 & 18 \\
10 & $[16, 19)$ & 18 & $16 < 18$ & 跳過 & \\
3 & $[17, 20)$ & 18 & $17 < 18$ & 跳過 & \\
12 & $[19, 21)$ & 18 & $19 \ge 18$ & 選 & 21 \\
\bottomrule
\end{tabular}
\end{center}

選到 $8, 6, 5, 11, 7, 12$ 共 $6$ 筆。 第 4 筆和第 5 筆都在 $13$ 點結束，依編號第 4 筆
排前面；它在 $8$ 點開始，和第 6 筆衝突，所以輪到它時被跳過。 $10$、$12$、$15$ 這三個時
刻，前一筆剛結束、下一筆立刻開始，在半開區間下不算衝突。

\begin{examnote}
區間題先看清楚端點：題目寫成閉區間 $[s_i, f_i]$，或寫「時刻 $f_i$ 還在使用」，接在後
面的申請就要 $s_j > f_i$ 才相容，虛擬碼~\ref{alg:gr-activity} 第 5 行的 $\ge$ 要改成
$>$，答案可能變少。 作答時寫一句你用的是哪一種，排序時結束時刻相同怎麼排也寫出來。
\end{examnote}

\subsection{為什麼最早結束的一定對}
\label{sec:gr-act-proof}

證明分兩步。 第一步說明「第一個選擇不會錯」，第二步說明「選了之後，剩下的是同一種
問題」。

\begin{theorem}
\label{thm:gr-act-choice}
設 $S$ 是一群申請，$a$ 是 $S$ 裡結束時刻最早的一筆（有好幾筆一樣早時任選一筆）。 那
麼 $S$ 有一個最大的相容子集合包含 $a$。
\end{theorem}

\begin{proof}
任取一個最大的相容子集合 $O$，令 $o$ 是 $O$ 裡結束最早的一筆。 如果 $o = a$，已經得
證。 否則令 $O' = (O - \{o\}) \cup \{a\}$，也就是把 $o$ 換成 $a$。

$O$ 裡其他的申請 $q$ 都和 $o$ 相容，而且 $f_q \ge f_o$。 $q$ 不可能整個在 $o$ 前面
（那樣 $f_q \le s_o < f_o$），所以 $q$ 在 $o$ 後面：$s_q \ge f_o$。 $a$ 是全部申請裡
結束最早的，$f_a \le f_o \le s_q$，所以 $a$ 和每一個 $q$ 都相容。 $O'$ 是相容子集合，
筆數和 $O$ 一樣，也是最大的，而且包含 $a$。
\end{proof}

\textbf{第二步。} 選了 $a$ 之後，還能選的是 $S' = \{q \in S : s_q \ge f_a\}$，其他申
請都和 $a$ 衝突。 如果 $O$ 是 $S$ 的最大相容子集合而且包含 $a$，那麼 $O - \{a\}$ 一
定是 $S'$ 的最大相容子集合：它是 $S'$ 的相容子集合；如果 $S'$ 有更大的相容子集合
$P$，那 $P \cup \{a\}$ 就比 $O$ 大，矛盾。

\textbf{合起來。} 對申請筆數做歸納。 虛擬碼~\ref{alg:gr-activity} 先選了 $S$ 裡結束
最早的 $a$；之後它跳過的申請都在 $f_a$ 之前開始，正好是不在 $S'$ 裡的那些，所以它接
下來做的事，就是對 $S'$ 從頭跑一次同一個演算法。 由歸納假設，它在 $S'$ 選到的筆數就是
$S'$ 的最大筆數 $m'$；由定理~\ref{thm:gr-act-choice} 和第二步，$S$ 的最大筆數是
$1 + m'$，貪心選到的也是 $1 + m'$。

\subsection{貪心法的兩個要素}
\label{sec:gr-elements}

上面的證明拆成兩件事，幾乎每個貪心演算法的證明都是這兩件：

\begin{enumerate}
\item \textbf{貪心選擇性質}（greedy-choice property）：存在一個最佳解，它的第一個選
擇就是貪心的選擇。 通常用交換論證：任取一個最佳解，把它和貪心不同的那個選擇換成貪心
的，說明換完仍然可行、而且不會變差。
\item \textbf{最佳子結構}（optimal substructure）：做了貪心的選擇之後，剩下的是一個同
類型、比較小的問題，而且原問題的最佳解扣掉那個選擇，正好是小問題的最佳解。
\end{enumerate}

第~18~章的動態規劃也需要最佳子結構。 兩者的差別在做選擇的時機：動態規劃先把所有子
問題的答案算出來，再比較各種選擇；貪心在解子問題之前就做了選擇，而且只留下一個子問
題。 所以貪心比較快，但前提是貪心選擇性質成立，這一點每個問題都要另外證明。

寫交換論證時常漏掉兩件事：換完之後解還是\textbf{可行}的（上面的證明花了一半篇幅說
$a$ 和其他申請都相容），以及第一步之後為什麼也對（要靠最佳子結構加歸納，或像觀念題
第~\ref{q:gr-ahead}~題那樣逐步比較）。

\subsection{C 程式}

\begin{cbookrecap}[第~14~章結構、第~18~章 18.6 節]
\texttt{qsort(base, n, size, cmp)} 排序從 \texttt{base} 開始的 \texttt{n} 個元素，每個
\texttt{size} 個位元組。 比較函式拿到兩個 \texttt{const void *}，先轉成指向元素型別的
指標再取欄位；第一個要排前面就回傳負數，後面就回傳正數，相等回傳 $0$。 用
\texttt{x->f - y->f} 回傳差值可能溢位（本書第~14~章 14.7 節的例題），所以下面的程式
一律用比較運算子。
\end{cbookrecap}

程式讀進 $n$ 筆申請，第 $i$ 行是第 $i$ 筆的 $s_i$ 和 $f_i$。 申請存成結構，排序時把編
號一起帶著走，邊掃描邊印出選中的申請。

\begin{example}[label=ex:gr-activity]{c}{(activity.c) 活動選擇：最早結束的先排}
#include <stdio.h>
#include <stdlib.h>

typedef struct {
    int id;         // 申請編號，從 1 開始
    int s, f;       // 借用時段 [s, f)
} Booking;

// 依結束時刻由小到大；結束時刻相同時編號小的在前。 不用減法，避免溢位。
int by_finish(const void *p, const void *q)
{
    const Booking *x = p, *y = q;
    if (x->f != y->f) {
        return x->f < y->f ? -1 : 1;
    }
    return (x->id > y->id) - (x->id < y->id);
}

int main(void)
{
    int n;
    if (scanf("%d", &n) != 1 || n < 0) {
        return 1;
    }
    Booking *b = malloc((n + 1) * sizeof *b);   // 多一格：n = 0 時不 malloc(0)
    if (b == NULL) {
        return 1;
    }
    for (int i = 0; i < n; i++) {
        if (scanf("%d %d", &b[i].s, &b[i].f) != 2 || b[i].s >= b[i].f) {
            free(b);
            return 1;
        }
        b[i].id = i + 1;
    }
    qsort(b, n, sizeof b[0], by_finish);
    int k = 0, last = 0;        // 選了幾筆、最後一筆選中的結束時刻
    for (int i = 0; i < n; i++) {
        if (k == 0 || b[i].s >= last) {         // 半開區間：接著開始不算衝突
            printf("#%d [%d, %d)\n", b[i].id, b[i].s, b[i].f);
            k++;
            last = b[i].f;
        }
    }
    printf("accepted %d of %d\n", k, n);
    free(b);
    return 0;
}
\end{example}
\iosample{12 \\ 13 16 \\ 9 11 \\ 17 20 \\ 8 13 \\ 12 13 \\ 10 12 \\ 15 18 \\ 8 10 \\ 11 14 \\ 16 19 \\ 14 15 \\ 19 21}{\#8 [8, 10) \\ \#6 [10, 12) \\ \#5 [12, 13) \\ \#11 [14, 15) \\ \#7 [15, 18) \\ \#12 [19, 21) \\ accepted 6 of 12}

輸出和手算的表一樣。 \texttt{malloc} 多要一格，是因為 C 標準允許 \texttt{malloc(0)}
回傳 \texttt{NULL}，$n = 0$ 時程式會誤以為要不到記憶體。

\section{分數背包與 0/1 背包}
\label{sec:gr-knapsack}

\subsection{分數背包}

咖啡豆批發行有 $n$ 種豆子，第 $i$ 種的庫存是 $w_i$ 公斤，整批的價值是 $v_i$ 元（$w_i$、
$v_i$ 都是正數）。 袋子最多裝 $W$ 公斤。 每一種豆子都可以只拿一部分：拿第 $i$ 種的比
例 $x_i$（$0 \le x_i \le 1$），得到 $x_i v_i$ 元，佔 $x_i w_i$ 公斤。 在
$\sum x_i w_i \le W$ 的條件下讓 $\sum x_i v_i$ 最大，這叫\textbf{分數背包問題}
（fractional knapsack problem）。

既然可以拆，每一公斤的價值才是重點。 令 $r_i = v_i / w_i$ 是第 $i$ 種的\textbf{單位
重量價值}。 貪心的做法是：依 $r_i$ 由大到小，能整批拿就整批拿，袋子快滿時拿最後一種
的一部分把袋子裝滿。

\begin{pseudo}[label=alg:gr-frac]{分數背包（已依 $v[i] / w[i]$ 由大到小排好）}
\Function{Fractional-Knapsack}{$w, v, n, W$}
    \State $\mathit{room} \gets W$
    \State $\mathit{total} \gets 0$
    \For{$i \gets 1$ \To $n$}
        \State $x_i \gets \min(1, \mathit{room} / w[i])$ \Comment{第 $i$ 種拿的比例}
        \State $\mathit{total} \gets \mathit{total} + x_i \cdot v[i]$
        \State $\mathit{room} \gets \mathit{room} - x_i \cdot w[i]$
    \EndFor
    \State \Return $\mathit{total}$
\EndFunction
\end{pseudo}

袋子滿了以後 $\mathit{room} = 0$，之後每一種的 $x_i$ 都是 $0$。 時間花在排序，
$\Theta(n \lg n)$。

\textbf{例子。} $W = 10$ 公斤，五種豆子：

\begin{center}
\begin{tabular}{cccccc}
\toprule
豆子 & A & B & C & D & E \\
\midrule
庫存 $w_i$（公斤） & 4 & 3 & 5 & 2 & 6 \\
整批價值 $v_i$（元） & 2000 & 1950 & 2250 & 1400 & 2100 \\
每公斤 $r_i$ & 500 & 650 & 450 & 700 & 350 \\
\bottomrule
\end{tabular}
\end{center}

依 $r_i$ 由大到小是 D、B、A、C、E。 D 整批拿（$2$ 公斤，$1400$ 元），B 整批拿（共
$5$ 公斤，$3350$ 元），A 整批拿（共 $9$ 公斤，$5350$ 元）。 剩 $1$ 公斤，拿 C 的
$1/5$，$450$ 元。 總共 $5800$ 元，E 完全沒拿。

\begin{theorem}
\label{thm:gr-frac}
虛擬碼~\ref{alg:gr-frac} 得到分數背包的最大總價值。
\end{theorem}

\begin{proof}
令 $x_1, \ldots, x_n$ 是貪心的解，$y_1, \ldots, y_n$ 是任何一個最佳解。 如果兩者不
同，取第一個不同的位置 $j$。 貪心在第 $j$ 種拿到「能拿的最多」，所以 $y_j < x_j$。
最佳解在第 $j$ 種少拿了 $\delta = (x_j - y_j) w_j$ 公斤。 前 $j - 1$ 種兩者一樣，貪心
拿完第 $j$ 種沒有超重，所以最佳解在第 $j + 1$ 種以後的重量加上袋子的空位，等於
$W - \sum_{i<j} x_i w_i - y_j w_j \ge x_j w_j - y_j w_j = \delta$。 先用袋子的空位，不夠的
部分再從第 $j + 1$ 種以後拿出豆子，湊滿 $\delta$ 公斤，換成第 $j$ 種的 $\delta$ 公斤。
被拿出的豆子每公斤價值都 $\le r_j$，空位的價值是 $0$，所以總價值不會變少，重量也沒有超
過。 換完之後仍是最佳解，而且和貪心解第一個不同的位置往後移了。 重複這個交換，最後最佳
解變成貪心解，所以貪心解也是最佳的。
\end{proof}

\subsection{0/1 背包}

把豆子換成不能拆的東西：每一件只能整件拿或不拿，$x_i \in \{0, 1\}$。 這是
\textbf{0/1 背包問題}（0/1 knapsack problem）。 照單位重量價值拿還對嗎？ 容量 $W = 8$，
三件物品：甲重 $5$、值 $60$（每單位 $12$），乙、丙都是重 $4$、值 $44$（每單位 $11$）。

依單位重量價值，先拿甲（$5$ 公斤，$60$），剩 $3$ 公斤，乙、丙都放不下，總價值 $60$。
不拿甲，拿乙和丙，剛好 $8$ 公斤，總價值 $88$。 貪心錯了。

「價值大的先拿」在這組資料上也是先拿甲；「輕的先拿」在容量 $8$、物品（重 $2$，值
$2$）和（重 $7$，值 $70$）上只得到 $2$。

問題出在定理~\ref{thm:gr-frac} 的證明哪一步？ 交換的時候要從別的物品「拿出 $\delta$
公斤」換成第 $j$ 種，這一步需要把東西切開。 0/1 背包不能切，拿了甲之後剩下的 $3$ 公
斤空間就浪費了，而這個浪費會不會發生，要看後面還有哪些物品，貪心在做第一個選擇時看
不到。 0/1 背包要用第~18~章的動態規劃，時間是 $O(nW)$，這是第~1~章 1.9 節說的擬多
項式時間；第~28~章會看到，一般的 0/1 背包是 NP 完備的，目前不知道有多項式時間的演算
法。

兩個問題都有最佳子結構，差別只在貪心選擇性質，可見「有最佳子結構」不足以讓貪心成立。

另一個用途：分數背包的最佳值一定不小於同一組資料的 0/1 背包最佳值（0/1 的解也是分數
的解），所以它可以當 0/1 背包的\textbf{上界}。 第~20~章的分支限界就用這個上界來剪
枝。

\section{霍夫曼編碼與最佳合併樹}
\label{sec:gr-huffman}

\subsection{定長碼與變長碼}

山上的氣象站每小時記錄一個天氣代碼，一共七種：晴 S、多雲 C、陰 O、雨 R、雷雨 T、霧
F、霾 H。 某一段時間的 $100$ 筆紀錄裡，各代碼出現的次數是：

\begin{center}
\begin{tabular}{c*{7}{c}}
\toprule
代碼 & S & C & O & R & T & F & H \\
\midrule
次數 & 34 & 22 & 15 & 13 & 8 & 5 & 3 \\
\bottomrule
\end{tabular}
\end{center}

要把紀錄存成 $0$、$1$ 組成的位元串。 七種代碼，每種給一個 $3$ 位元的碼（$2^3 = 8 \ge
7$）就夠了，$100$ 筆共 $300$ 位元。 這叫\textbf{定長碼}（fixed-length code）。 晴天
佔了三分之一，如果給 S 比較短的碼、給少見的 H 比較長的碼，總長度可以更短，這叫
\textbf{變長碼}（variable-length code）。

變長碼要能解得開。 如果 S 的碼是 \texttt{1}、C 的碼是 \texttt{10}，讀到 \texttt{10}
時不知道是 C，還是 S 後面接著另一個碼。 最簡單的保證是：\textbf{沒有一個碼是另一個碼
的前綴}（開頭的一段）。 這樣的碼叫\textbf{前綴碼}（prefix code）。 解前綴碼時從頭讀，
一讀到某個完整的碼就可以確定是它，不必往後看。 任何能唯一解碼的碼，都有碼長完全一樣的前綴碼（Kraft--McMillan 不等式），所以只看前綴
碼不會吃虧。

\subsection{碼樹}

前綴碼可以畫成一棵二元樹：從根出發，往左走記 \texttt{0}，往右走記 \texttt{1}，每個
符號放在一個葉節點，根到葉節點的路徑就是它的碼。 符號都在葉節點上，所以沒有一個碼是
另一個的前綴；反過來，任何前綴碼都能這樣畫成一棵樹。 解碼就是從根往下走，走到葉節點
就輸出那個符號，再回到根。

\textbf{碼長和深度。} 本書照第~7~章的定義：節點的\textbf{深度}數的是從根往下走的
\textbf{邊數}，根的深度是 $0$；樹的\textbf{高度}數的是\textbf{節點數}（層數）。 每走
一條邊寫一個位元，所以
\[
\text{符號 $c$ 的碼長} = \text{葉節點 $c$ 的深度 } d_T(c) \text{（邊數）}, \qquad
\text{最長的碼長} = \text{樹的高度} - 1.
\]
符號 $c$ 出現 $c.\mathit{freq}$ 次，整份資料編碼後的總位元數是
\[
B(T) = \sum_{c \in C} c.\mathit{freq} \cdot d_T(c),
\]
這也叫這棵樹的\textbf{加權路徑長}（weighted path length，只算葉節點）。 我們要找
$B(T)$ 最小的樹，對應的碼叫\textbf{最佳前綴碼}。

\textbf{最佳的碼樹是嚴格二元樹。} 如果某個內部節點只有一個孩子，把它刪掉、讓孩子直接
接上去，它底下每個葉節點的深度都少 $1$，$B(T)$ 變小（頻率都是正的）。 所以最佳的碼
樹每個內部節點都有兩個孩子，是第~7~章的嚴格二元樹；$n$ 個符號的碼樹有 $n$ 個葉節點、
$n - 1$ 個內部節點（$n_0 = n_2 + 1$）。 定長碼的樹一般不是嚴格的：七個符號用 $3$ 位
元，有一個碼（例如 \texttt{111}）沒有人用，它的兄弟可以往上提一層。

\subsection{霍夫曼演算法}

霍夫曼（D. A. Huffman，1952）的做法由下往上建樹。 一開始每個符號自己是一棵只有一個
節點的樹，權重是它的頻率。 每一步拿出權重最小的兩棵樹，接到一個新的根底下，新根的權
重是兩者的和，放回去。 每一步樹少一棵，$n - 1$ 步之後只剩一棵，就是碼樹。

「拿出最小的兩個、放回一個」正是第~8~章最小優先佇列的操作。 權重一樣的時候要有固定
的規則，否則同一組資料可以建出不同的樹。 本書的規則：

\begin{itemize}
\item 每個節點有一個編號 $\mathit{id}$。 葉節點照輸入的順序編 $0, 1, \ldots, n - 1$，
內部節點照建立的順序接著編 $n, n + 1, \ldots$。
\item 優先佇列先比權重，權重一樣時比編號，編號小的先出來。 所以權重相同時，葉節點比內
部節點先出來，早建的內部節點比晚建的先出來。
\item 先取出的當左孩子（邊標 \texttt{0}），後取出的當右孩子（邊標 \texttt{1}）。
\end{itemize}

\begin{pseudo}[label=alg:gr-huffman]{霍夫曼編碼}
\Function{Huffman}{$C$} \Comment{$C$：$n \ge 2$ 個符號，各有 $\mathit{freq}$}
    \State 建立最小優先佇列 $Q$，放進 $C$ 的每個符號（鍵是 $(\mathit{freq}, \mathit{id})$）
    \State $k \gets |C|$ \Comment{下一個內部節點的編號}
    \While{$Q$ 裡不只一個節點}
        \State $x \gets \Call{Extract-Min}{Q}$
        \State $y \gets \Call{Extract-Min}{Q}$
        \State 建立節點 $z$：$z.\mathit{left} \gets x$，$z.\mathit{right} \gets y$，
        $z.\mathit{freq} \gets x.\mathit{freq} + y.\mathit{freq}$，$z.\mathit{id} \gets k$
        \State $k \gets k + 1$
        \State \Call{Insert}{$Q, z$}
    \EndWhile
    \State \Return \Call{Extract-Min}{$Q$} \Comment{碼樹的根}
\EndFunction
\end{pseudo}

用二元堆積實作，第 2 行由下而上建堆積是 $O(n)$（第~8~章 8.2 節），迴圈跑 $n - 1$ 次，
每次三個堆積操作，各 $O(\lg n)$，總共 $O(n \lg n)$。

\subsection{逐步建樹}
\label{sec:gr-huff-trace}

用氣象站的資料跑一次。 下面每一張圖是一次合併之後，優先佇列裡的所有樹，由左到右照
$(\mathit{freq}, \mathit{id})$ 排列，所以下一步要取出的永遠是最左邊兩棵。 方框是葉節
點（代碼和次數），圓圈是內部節點（權重），橘色是這一步新建的節點。 圖是用程式照上面
的規則模擬後畫出來的。

\begin{center}
\begin{minipage}{\linewidth}
\centering
\begin{tikzpicture}[baseline=(current bounding box.north)]
\node[gr-leaf] (ha0) at (6.600,0.000) {S\\34};
\node[gr-leaf] (ha1) at (5.500,0.000) {C\\22};
\node[gr-leaf] (ha2) at (4.400,0.000) {O\\15};
\node[gr-leaf] (ha3) at (3.300,0.000) {R\\13};
\node[gr-leaf] (ha4) at (2.200,0.000) {T\\8};
\node[gr-leaf] (ha5) at (1.100,0.000) {F\\5};
\node[gr-leaf] (ha6) at (0.000,0.000) {H\\3};
\node[gr-plab] at (3.300,-1.750) {(a) 開始：七棵單節點的樹};
\end{tikzpicture}\hspace{0.7cm}\begin{tikzpicture}[baseline=(current bounding box.north)]
\node[gr-leaf] (hb0) at (6.320,0.000) {S\\34};
\node[gr-leaf] (hb1) at (5.220,0.000) {C\\22};
\node[gr-leaf] (hb2) at (4.120,0.000) {O\\15};
\node[gr-leaf] (hb3) at (3.020,0.000) {R\\13};
\node[gr-leaf] (hb4) at (0.000,0.000) {T\\8};
\node[gr-leaf] (hb5) at (1.920,-0.900) {F\\5};
\node[gr-leaf] (hb6) at (1.100,-0.900) {H\\3};
\node[gr-inner,gr-new] (hb7) at (1.510,0.000) {8};
\draw[gr-edge] (hb7) -- node[gr-bit, left] {0} (hb6);
\draw[gr-edge] (hb7) -- node[gr-bit, right] {1} (hb5);
\node[gr-plab] at (3.160,-1.750) {(b) 第 1 次：$3 + 5 = 8$};
\end{tikzpicture}

\bigskip
\begin{tikzpicture}[baseline=(current bounding box.north)]
\node[gr-leaf] (hc0) at (6.040,0.000) {S\\34};
\node[gr-leaf] (hc1) at (4.940,0.000) {C\\22};
\node[gr-leaf] (hc2) at (1.100,0.000) {O\\15};
\node[gr-leaf] (hc3) at (0.000,0.000) {R\\13};
\node[gr-leaf] (hc4) at (2.200,-0.900) {T\\8};
\node[gr-leaf] (hc5) at (3.840,-1.800) {F\\5};
\node[gr-leaf] (hc6) at (3.020,-1.800) {H\\3};
\node[gr-inner] (hc7) at (3.430,-0.900) {8};
\node[gr-inner,gr-new] (hc8) at (2.815,0.000) {16};
\draw[gr-edge] (hc7) -- node[gr-bit, left] {0} (hc6);
\draw[gr-edge] (hc7) -- node[gr-bit, right] {1} (hc5);
\draw[gr-edge] (hc8) -- node[gr-bit, left] {0} (hc4);
\draw[gr-edge] (hc8) -- node[gr-bit, right] {1} (hc7);
\node[gr-plab] at (3.020,-2.650) {(c) 第 2 次：$8 + 8 = 16$};
\end{tikzpicture}\hspace{0.7cm}\begin{tikzpicture}[baseline=(current bounding box.north)]
\node[gr-leaf] (hd0) at (5.760,0.000) {S\\34};
\node[gr-leaf] (hd1) at (2.740,0.000) {C\\22};
\node[gr-leaf] (hd2) at (4.660,-0.900) {O\\15};
\node[gr-leaf] (hd3) at (3.840,-0.900) {R\\13};
\node[gr-leaf] (hd4) at (0.000,-0.900) {T\\8};
\node[gr-leaf] (hd5) at (1.640,-1.800) {F\\5};
\node[gr-leaf] (hd6) at (0.820,-1.800) {H\\3};
\node[gr-inner] (hd7) at (1.230,-0.900) {8};
\node[gr-inner] (hd8) at (0.615,0.000) {16};
\node[gr-inner,gr-new] (hd9) at (4.250,0.000) {28};
\draw[gr-edge] (hd7) -- node[gr-bit, left] {0} (hd6);
\draw[gr-edge] (hd7) -- node[gr-bit, right] {1} (hd5);
\draw[gr-edge] (hd8) -- node[gr-bit, left] {0} (hd4);
\draw[gr-edge] (hd8) -- node[gr-bit, right] {1} (hd7);
\draw[gr-edge] (hd9) -- node[gr-bit, left] {0} (hd3);
\draw[gr-edge] (hd9) -- node[gr-bit, right] {1} (hd2);
\node[gr-plab] at (2.880,-2.650) {(d) 第 3 次：$13 + 15 = 28$};
\end{tikzpicture}
\diagcaption{霍夫曼演算法 (a)--(d)：開始和前三次合併。 每一張圖由左到右是優先佇列裡的樹，最左邊兩棵就是下一步要取出的。 邊上的 \texttt{0}、\texttt{1} 是先取出（左）和後取出（右）。}
\label{fig:gr-huff1}
\end{minipage}
\end{center}

第 2 次合併是平手的例子：佇列裡有兩個權重 $8$ 的節點，葉節點 T（編號 $4$）和第 1 次
建的內部節點（編號 $7$）。 照規則葉節點先出來，T 當左孩子。

\begin{center}
\begin{minipage}{\linewidth}
\centering
\begin{tikzpicture}[baseline=(current bounding box.north)]
\node[gr-leaf] (he0) at (1.920,0.000) {S\\34};
\node[gr-leaf] (he1) at (5.480,-0.900) {C\\22};
\node[gr-leaf] (he2) at (0.820,-0.900) {O\\15};
\node[gr-leaf] (he3) at (0.000,-0.900) {R\\13};
\node[gr-leaf] (he4) at (3.020,-1.800) {T\\8};
\node[gr-leaf] (he5) at (4.660,-2.700) {F\\5};
\node[gr-leaf] (he6) at (3.840,-2.700) {H\\3};
\node[gr-inner] (he7) at (4.250,-1.800) {8};
\node[gr-inner] (he8) at (3.635,-0.900) {16};
\node[gr-inner] (he9) at (0.410,0.000) {28};
\node[gr-inner,gr-new] (he10) at (4.558,0.000) {38};
\draw[gr-edge] (he9) -- node[gr-bit, left] {0} (he3);
\draw[gr-edge] (he9) -- node[gr-bit, right] {1} (he2);
\draw[gr-edge] (he7) -- node[gr-bit, left] {0} (he6);
\draw[gr-edge] (he7) -- node[gr-bit, right] {1} (he5);
\draw[gr-edge] (he8) -- node[gr-bit, left] {0} (he4);
\draw[gr-edge] (he8) -- node[gr-bit, right] {1} (he7);
\draw[gr-edge] (he10) -- node[gr-bit, left] {0} (he8);
\draw[gr-edge] (he10) -- node[gr-bit, right] {1} (he1);
\node[gr-plab] at (2.740,-3.550) {(e) 第 4 次：$16 + 22 = 38$};
\end{tikzpicture}\hspace{0.7cm}\begin{tikzpicture}[baseline=(current bounding box.north)]
\node[gr-leaf] (hf0) at (5.200,-0.900) {S\\34};
\node[gr-leaf] (hf1) at (2.460,-0.900) {C\\22};
\node[gr-leaf] (hf2) at (4.380,-1.800) {O\\15};
\node[gr-leaf] (hf3) at (3.560,-1.800) {R\\13};
\node[gr-leaf] (hf4) at (0.000,-1.800) {T\\8};
\node[gr-leaf] (hf5) at (1.640,-2.700) {F\\5};
\node[gr-leaf] (hf6) at (0.820,-2.700) {H\\3};
\node[gr-inner] (hf7) at (1.230,-1.800) {8};
\node[gr-inner] (hf8) at (0.615,-0.900) {16};
\node[gr-inner] (hf9) at (3.970,-0.900) {28};
\node[gr-inner] (hf10) at (1.538,0.000) {38};
\node[gr-inner,gr-new] (hf11) at (4.585,0.000) {62};
\draw[gr-edge] (hf7) -- node[gr-bit, left] {0} (hf6);
\draw[gr-edge] (hf7) -- node[gr-bit, right] {1} (hf5);
\draw[gr-edge] (hf8) -- node[gr-bit, left] {0} (hf4);
\draw[gr-edge] (hf8) -- node[gr-bit, right] {1} (hf7);
\draw[gr-edge] (hf10) -- node[gr-bit, left] {0} (hf8);
\draw[gr-edge] (hf10) -- node[gr-bit, right] {1} (hf1);
\draw[gr-edge] (hf9) -- node[gr-bit, left] {0} (hf3);
\draw[gr-edge] (hf9) -- node[gr-bit, right] {1} (hf2);
\draw[gr-edge] (hf11) -- node[gr-bit, left] {0} (hf9);
\draw[gr-edge] (hf11) -- node[gr-bit, right] {1} (hf0);
\node[gr-plab] at (2.600,-3.550) {(f) 第 5 次：$28 + 34 = 62$};
\end{tikzpicture}

\bigskip
\begin{tikzpicture}[baseline=(current bounding box.north)]
\node[gr-leaf] (hg0) at (4.920,-1.800) {S\\34};
\node[gr-leaf] (hg1) at (2.460,-1.800) {C\\22};
\node[gr-leaf] (hg2) at (4.100,-2.700) {O\\15};
\node[gr-leaf] (hg3) at (3.280,-2.700) {R\\13};
\node[gr-leaf] (hg4) at (0.000,-2.700) {T\\8};
\node[gr-leaf] (hg5) at (1.640,-3.600) {F\\5};
\node[gr-leaf] (hg6) at (0.820,-3.600) {H\\3};
\node[gr-inner] (hg7) at (1.230,-2.700) {8};
\node[gr-inner] (hg8) at (0.615,-1.800) {16};
\node[gr-inner] (hg9) at (3.690,-1.800) {28};
\node[gr-inner] (hg10) at (1.538,-0.900) {38};
\node[gr-inner] (hg11) at (4.305,-0.900) {62};
\node[gr-inner,gr-new] (hg12) at (2.921,0.000) {100};
\draw[gr-edge] (hg7) -- node[gr-bit, left] {0} (hg6);
\draw[gr-edge] (hg7) -- node[gr-bit, right] {1} (hg5);
\draw[gr-edge] (hg8) -- node[gr-bit, left] {0} (hg4);
\draw[gr-edge] (hg8) -- node[gr-bit, right] {1} (hg7);
\draw[gr-edge] (hg10) -- node[gr-bit, left] {0} (hg8);
\draw[gr-edge] (hg10) -- node[gr-bit, right] {1} (hg1);
\draw[gr-edge] (hg9) -- node[gr-bit, left] {0} (hg3);
\draw[gr-edge] (hg9) -- node[gr-bit, right] {1} (hg2);
\draw[gr-edge] (hg11) -- node[gr-bit, left] {0} (hg9);
\draw[gr-edge] (hg11) -- node[gr-bit, right] {1} (hg0);
\draw[gr-edge] (hg12) -- node[gr-bit, left] {0} (hg10);
\draw[gr-edge] (hg12) -- node[gr-bit, right] {1} (hg11);
\node[gr-plab] at (2.460,-4.450) {(g) 第 6 次：$38 + 62 = 100$};
\end{tikzpicture}
\diagcaption{霍夫曼演算法 (e)--(g)：第 4 到第 6 次合併。 (g) 是最後的碼樹，根的權重 $100$ 是紀錄的總筆數。 從根走到葉節點，把邊上的位元串起來就是那個代碼的碼。}
\label{fig:gr-huff2}
\end{minipage}
\end{center}
第 4 次合併時，權重 $16$ 的內部節點比 C 的 $22$ 小，所以它當左孩子，C 當右孩子；所以
C 的碼是 \texttt{01}，不是 \texttt{00}。

從圖 (g) 讀出每個代碼的碼，再和定長碼比較：

\begin{center}
\begin{tabular}{cccccc}
\toprule
代碼 & 次數 & 霍夫曼碼 & 碼長（深度） & 位元數 & 定長碼的位元數 \\
\midrule
S & 34 & \texttt{11} & 2 & 68 & 102 \\
C & 22 & \texttt{01} & 2 & 44 & 66 \\
O & 15 & \texttt{101} & 3 & 45 & 45 \\
R & 13 & \texttt{100} & 3 & 39 & 39 \\
T & 8 & \texttt{000} & 3 & 24 & 24 \\
F & 5 & \texttt{0011} & 4 & 20 & 15 \\
H & 3 & \texttt{0010} & 4 & 12 & 9 \\
\midrule
合計 & 100 & & & 252 & 300 \\
\bottomrule
\end{tabular}
\end{center}

總共 $252$ 位元，平均每筆 $2.52$ 位元，比定長碼省了 $16\%$。 F 和 H 的碼反而比定長碼
長，但它們很少出現。 驗算 $B(T)$ 有另一個方法：它等於所有\textbf{內部節點}的權重和，
$8 + 16 + 28 + 38 + 62 + 100 = 252$。 原因是每個符號被它的每一個祖先（內部節點）各算
一次，祖先的個數正好是它的深度；第~\ref{sec:gr-merge}~節會再用到這個事實。

\textbf{解碼。} 位元串 \texttt{1101100101110010} 從根開始走：\texttt{11} 到 S，回到
根；\texttt{01} 到 C；\texttt{100} 到 R；\texttt{101} 到 O；\texttt{11} 到 S；
\texttt{0010} 到 H。 解出 SCROSH。

\begin{examnote}
霍夫曼樹的追蹤題，權重平手時選哪一個、誰當左孩子，會改變樹的形狀和每個符號的碼，甚至
改變碼長。 例如頻率 $1, 1, 2, 2$：先合併兩個 $1$ 得到 $2$ 以後，三個 $2$ 平手。 合併
兩個原本的 $2$，四個碼長都是 $2$；合併新的 $2$ 和一個原本的 $2$，碼長是 $1, 2, 3, 3$。
兩種樹的總位元數都是 $12$，都是最佳的。 不同課本的平手規則和左右規則不同，作答時照題
目的規定；題目沒有規定，就先寫出你用的規則（例如本書的「先比權重再比編號，先取出的放
左邊」），閱卷的人才能照你的規則核對。 總位元數不受這些規則影響，可以拿來驗算。
\end{examnote}

\subsection{為什麼霍夫曼碼最佳}
\label{sec:gr-huff-proof}

證明的架構和活動選擇一樣：第一個引理是貪心選擇性質，第二個是最佳子結構。 以下 $C$
是符號的集合，頻率都是正數，$|C| \ge 2$。

\begin{lemma}
\label{lem:gr-huff-choice}
設 $x$、$y$ 是 $C$ 裡頻率最小的兩個符號。 存在一棵最佳碼樹，$x$ 和 $y$ 是兄弟，而且
都在最深的那一層。
\end{lemma}

\begin{proof}
任取一棵最佳碼樹 $T$。 它是嚴格二元樹，所以最深的葉節點 $a$ 一定有兄弟 $b$；$b$ 也是
葉節點，否則 $b$ 底下的葉節點比 $a$ 更深。 $d_T(a) = d_T(b)$ 是最大的深度。

想法是把 $x$、$y$ 搬到 $a$、$b$ 的位置。 先看交換兩個葉節點 $u$、$v$ 的效果：只有它們
的深度改變，
\[
B(\text{換後}) - B(\text{換前}) = \bigl(u.\mathit{freq} - v.\mathit{freq}\bigr)
\bigl(d(v) - d(u)\bigr),
\]
所以把頻率比較小的符號換到比較深（或一樣深）的位置，$B$ 不會變大。

如果 $x$ 已經是 $a$ 或 $b$ 其中一個，把名字對調讓 $x = a$；$y$ 也一樣，讓 $y = b$。
剩下沒對好的，逐一交換：$x \ne a$ 時交換 $x$ 和 $a$，$y \ne b$ 時交換 $y$ 和 $b$。 名字
對調過之後，被換上去的 $a$ 或 $b$ 都不是 $x$、$y$，頻率不小於 $y.\mathit{freq} \ge
x.\mathit{freq}$；被換下去的 $x$ 或 $y$ 到了最深的位置。 由上面的式子，每次交換 $B$ 都
不會變大。 換完的樹仍是最佳的，$x$、$y$ 是兄弟，都在最深的一層。
\end{proof}

\begin{lemma}
\label{lem:gr-huff-sub}
設 $x$、$y$ 是 $C$ 裡頻率最小的兩個符號。 令 $C'$ 是把 $x$、$y$ 拿掉、加入一個新符號
$z$，$z.\mathit{freq} = x.\mathit{freq} + y.\mathit{freq}$。 如果 $T'$ 是 $C'$ 的
最佳碼樹，把 $T'$ 的葉節點 $z$ 換成一個內部節點、底下接 $x$ 和 $y$，得到的樹 $T$ 是
$C$ 的最佳碼樹。
\end{lemma}

\begin{proof}
先比較 $T$ 和 $T'$。 除了 $x$、$y$、$z$，每個符號的深度都一樣；$x$、$y$ 的深度是
$d_{T'}(z) + 1$。 所以
\[
B(T) = B(T') - z.\mathit{freq} \cdot d_{T'}(z)
+ (x.\mathit{freq} + y.\mathit{freq})\bigl(d_{T'}(z) + 1\bigr)
= B(T') + x.\mathit{freq} + y.\mathit{freq}.
\]
同樣的計算對任何一棵「$x$、$y$ 是兄弟」的碼樹都成立：把那對兄弟收回成一個葉節點
$z$，$B$ 少 $x.\mathit{freq} + y.\mathit{freq}$。

假設 $T$ 不是最佳的。 由引理~\ref{lem:gr-huff-choice}，有一棵最佳碼樹 $T^*$ 讓 $x$、$y$
是兄弟，而且 $B(T^*) < B(T)$。 把 $T^*$ 的 $x$、$y$ 收回成 $z$，得到 $C'$ 的碼樹
$T^{*\prime}$，
\[
B(T^{*\prime}) = B(T^*) - x.\mathit{freq} - y.\mathit{freq}
< B(T) - x.\mathit{freq} - y.\mathit{freq} = B(T'),
\]
和 $T'$ 是最佳的矛盾。
\end{proof}

\begin{theorem}
\label{thm:gr-huffman}
虛擬碼~\ref{alg:gr-huffman} 建出的碼樹是最佳的。
\end{theorem}

\begin{proof}
對符號個數 $n$ 做歸納。 $n = 2$ 時只有一棵嚴格二元樹（根接兩個葉節點），是最佳的。
$n > 2$ 時，演算法第一次合併取出頻率最小的兩個 $x$、$y$，建出 $z$；之後的每一步，和
「對 $C'$ 從頭跑霍夫曼演算法」做的事完全一樣（$z$ 當成一個葉節點）。 由歸納假設，它
對 $C'$ 建出最佳碼樹 $T'$；最後的樹就是 $T'$ 的 $z$ 展開成 $x$、$y$，由
引理~\ref{lem:gr-huff-sub}，它是 $C$ 的最佳碼樹。
\end{proof}

證明裡沒有用到平手規則：不管平手時怎麼選，取出的都是「頻率最小的兩個」，兩個引理照樣
成立。 這就是上面作答提醒說的，不同規則建出的樹都是最佳的。

\subsection{C 程式}

\begin{cbookrecap}[第~14~章結構、第~15~章遞迴、第~19~章 19.3 節]
\texttt{struct Node} 裡有兩個指向 \texttt{struct Node} 的指標，是建二元樹的寫法。 這裡一次要
一個 $2n - 1$ 格的陣列 \texttt{pool}，不替每個節點各 \texttt{malloc} 一次，前 $n$ 格
放葉節點、後 $n - 1$ 格放內部節點，節點之間仍用指標連起來，用完 \texttt{free} 一次就好。
\texttt{collect} 是遞迴的前序走訪。
\end{cbookrecap}

堆積照第~8~章的陣列堆積寫，改成最小堆積，存的是節點的指標，比較時用
\texttt{before}：先比頻率，再比編號。 碼存在一個 $n \times n$ 的字元陣列，第 $i$ 列
是編號 $i$ 的葉節點的碼；碼長最多 $n - 1$，加上結尾的 \verb|'\0'| 剛好 $n$ 格。 印的
時候照輸入的順序，結果固定。 輸入是符號個數 $n \ge 2$，以及每個符號和它的頻率。

\begin{example}[label=ex:gr-huffman]{c}{(huffman.c) 用最小堆積建霍夫曼樹，照輸入順序印出每個碼}
#include <stdio.h>
#include <stdlib.h>
#include <string.h>

typedef struct Node {
    long long w;                // 頻率（內部節點是兩個孩子的和）
    int id;                     // 葉節點 0..n-1 照輸入順序，內部節點照建立順序
    char sym;                   // 只有葉節點用到
    struct Node *left, *right;  // 葉節點兩個都是 NULL
} Node;

// 平手規則：先比頻率，頻率相同比 id。 id 都不同，所以建出的樹是唯一的。
static int before(const Node *x, const Node *y)
{
    return x->w < y->w || (x->w == y->w && x->id < y->id);
}

// 最小堆積 h[0..*size-1] 存節點的指標。 上浮、下沉同第 8 章，
// 只是把要放的節點先拿在手上，一路把父（子）節點挪過來，最後才放進去。
void push(Node **h, int *size, Node *x)
{
    int i = (*size)++;
    while (i > 0 && before(x, h[(i - 1) / 2])) {
        h[i] = h[(i - 1) / 2];
        i = (i - 1) / 2;
    }
    h[i] = x;
}

Node *pop(Node **h, int *size)          // 呼叫前要確定堆積不是空的
{
    Node *top = h[0], *x = h[--(*size)];
    int i = 0, j;
    while ((j = 2 * i + 1) < *size) {
        if (j + 1 < *size && before(h[j + 1], h[j])) {
            j++;                        // 取較小的孩子
        }
        if (!before(h[j], x)) {
            break;
        }
        h[i] = h[j];
        i = j;
    }
    h[i] = x;
    return top;
}

// 前序走訪，buf[0..depth-1] 是從根到 t 的路徑；到了葉節點就抄進 code 的第 id 列。
void collect(const Node *t, char *buf, int depth, char *code, int n)
{
    if (t->left == NULL) {
        buf[depth] = '\0';
        memcpy(code + (size_t)t->id * n, buf, depth + 1);
        return;
    }
    buf[depth] = '0';
    collect(t->left, buf, depth + 1, code, n);
    buf[depth] = '1';
    collect(t->right, buf, depth + 1, code, n);
}

int main(void)
{
    int n, size = 0;
    if (scanf("%d", &n) != 1 || n < 2) {
        return 1;
    }
    Node *pool = malloc((2 * n - 1) * sizeof *pool);   // n 個葉、n - 1 個內部節點
    Node **heap = malloc(n * sizeof *heap);
    char *code = malloc((size_t)n * n);    // 碼最長 n - 1 位，加 '\0' 共 n 格
    char *buf = malloc(n);
    int ok = pool != NULL && heap != NULL && code != NULL && buf != NULL;
    for (int i = 0; ok && i < n; i++) {
        ok = scanf(" %c %lld", &pool[i].sym, &pool[i].w) == 2 && pool[i].w > 0;
        pool[i].id = i;
        pool[i].left = pool[i].right = NULL;
        if (ok) {
            push(heap, &size, &pool[i]);
        }
    }
    if (ok) {
        for (int k = n; k < 2 * n - 1; k++) {           // n - 1 次合併
            Node *z = &pool[k];
            z->left = pop(heap, &size);                 // 先取出的當左孩子（邊標 0）
            z->right = pop(heap, &size);
            z->w = z->left->w + z->right->w;
            z->id = k;
            push(heap, &size, z);
            printf("merge %d: %lld + %lld = %lld\n",
                   k - n + 1, z->left->w, z->right->w, z->w);
        }
        Node *root = pop(heap, &size);
        collect(root, buf, 0, code, n);
        long long total = 0;
        for (int i = 0; i < n; i++) {                   // 照輸入順序印，結果固定
            const char *c = code + (size_t)i * n;
            printf("%c %lld %s\n", pool[i].sym, pool[i].w, c);
            total += pool[i].w * (long long)strlen(c);
        }
        printf("total bits = %lld\n", total);
    }
    free(pool);
    free(heap);
    free(code);
    free(buf);
    return ok ? 0 : 1;
}
\end{example}
\iosample{7 \\ S 34 \\ C 22 \\ O 15 \\ R 13 \\ T 8 \\ F 5 \\ H 3}{merge 1: 3 + 5 = 8 \\ merge 2: 8 + 8 = 16 \\ merge 3: 13 + 15 = 28 \\ merge 4: 16 + 22 = 38 \\ merge 5: 28 + 34 = 62 \\ merge 6: 38 + 62 = 100 \\ S 34 11 \\ C 22 01 \\ O 15 101 \\ R 13 100 \\ T 8 000 \\ F 5 0011 \\ H 3 0010 \\ total bits = 252}

六次合併和圖 (b)--(g) 一致，每個碼和上面的表相同。 堆積最多放 $n$ 個指標（每次合併
取出兩個、放回一個），所以 \texttt{heap} 只要 $n$ 格。 只有一個符號時碼樹只有根，碼長
是 $0$，沒辦法編碼，所以程式要求 $n \ge 2$。

頻率已經由小到大排好時可以不用堆積：葉節點放一個佇列，新建的內部節點依序放第二個佇
列（它們的權重自動遞增），每次比較兩個佇列的前端，總共 $O(n)$。

\subsection{最佳合併樹}
\label{sec:gr-merge}

第~15~章 15.3 節的外部排序先產生一些排好的串，再一次合併 $k$ 個，直到剩一個。 那裡
假設每個串一樣長，只需要算趟數；用取代選擇產生的串長短不一，「先合併哪幾個」就會影響
總成本。 第~15~章最後留了一句：這和霍夫曼編碼是同一個問題。 這一小節把它說清楚。

先看 $2$ 路合併。 有 $r$ 個排好的串，長度（紀錄筆數）是 $\ell_1, \ldots, \ell_r$。 合
併兩個長度 $a$、$b$ 的串，每筆紀錄讀一次、寫一次，成本算 $a + b$。 每次挑兩個串合併，
$r - 1$ 次之後剩一個。 整個合併過程可以畫成一棵\textbf{合併樹}：每個初始串是一個葉節點，
每次合併是一個內部節點，權重是合併出來的串長；這是一棵有 $r$ 個葉節點的嚴格二元樹。
總成本是所有內部節點的權重和。

\begin{lemma}
\label{lem:gr-merge-cost}
合併樹 $T$ 的總合併成本等於 $\sum_{i=1}^{r} \ell_i \cdot d_T(i)$，$d_T(i)$ 是第 $i$
個串的葉節點的深度（邊數）。
\end{lemma}

\begin{proof}
總成本是每次合併搬動的紀錄數加起來。 換個方向數：第 $i$ 個串的每一筆紀錄，在它的每
一個祖先（內部節點）那一次合併各被搬一次，祖先的個數就是它的深度 $d_T(i)$。 所以總成
本 $= \sum_i \ell_i \, d_T(i)$。
\end{proof}

右邊這個和，把 $\ell_i$ 看成頻率，正是上一小節的 $B(T)$，也叫這棵樹的\textbf{加權
外部路徑長}（weighted external path length）：外部節點（葉節點）的深度乘上權重的和。
所以「怎麼兩兩合併最省」和「怎麼建最佳碼樹」是同一個問題，答案也一樣：\textbf{每次合
併目前最短的兩個串}。 這樣合併出來的樹叫\textbf{最佳合併樹}（optimal merge tree）。

\textbf{例子。} 取代選擇產生了七個串，長度（單位：千筆）如下：

\begin{center}
\begin{tabular}{c*{7}{c}}
\toprule
串 & R1 & R2 & R3 & R4 & R5 & R6 & R7 \\
\midrule
長度 & 12 & 30 & 8 & 25 & 5 & 18 & 40 \\
\bottomrule
\end{tabular}
\end{center}

每次合併最短的兩個（一樣長時初始串先取，和霍夫曼演算法相同）：$5 + 8 = 13$，
$12 + 13 = 25$，$18 + 25 = 43$（R6 和 R4；兩個 $25$ 平手，先取初始串 R4），
$25 + 30 = 55$，$40 + 43 = 83$，$55 + 83 = 138$。 總成本 $357$ 千筆。 畫成合併樹：

\begin{center}
\begin{minipage}{\linewidth}
\centering
\begin{tikzpicture}
\node[gr-leaf] (mt0) at (0.000,-2.700) {R1\\12};
\node[gr-leaf] (mt1) at (2.850,-1.800) {R2\\30};
\node[gr-leaf] (mt2) at (1.900,-3.600) {R3\\8};
\node[gr-leaf] (mt3) at (5.700,-2.700) {R4\\25};
\node[gr-leaf] (mt4) at (0.950,-3.600) {R5\\5};
\node[gr-leaf] (mt5) at (4.750,-2.700) {R6\\18};
\node[gr-leaf] (mt6) at (3.800,-1.800) {R7\\40};
\node[gr-inner] (mt7) at (1.425,-2.700) {13};
\node[gr-inner] (mt8) at (0.712,-1.800) {25};
\node[gr-inner] (mt9) at (5.225,-1.800) {43};
\node[gr-inner] (mt10) at (1.781,-0.900) {55};
\node[gr-inner] (mt11) at (4.512,-0.900) {83};
\node[gr-inner] (mt12) at (3.147,0.000) {138};
\draw[gr-edge] (mt7) -- (mt4);
\draw[gr-edge] (mt7) -- (mt2);
\draw[gr-edge] (mt8) -- (mt0);
\draw[gr-edge] (mt8) -- (mt7);
\draw[gr-edge] (mt10) -- (mt8);
\draw[gr-edge] (mt10) -- (mt1);
\draw[gr-edge] (mt9) -- (mt5);
\draw[gr-edge] (mt9) -- (mt3);
\draw[gr-edge] (mt11) -- (mt6);
\draw[gr-edge] (mt11) -- (mt9);
\draw[gr-edge] (mt12) -- (mt10);
\draw[gr-edge] (mt12) -- (mt11);
\node[idxlbl, anchor=east] at (-0.6,0.000) {深度 0};
\node[idxlbl, anchor=east] at (-0.6,-0.900) {深度 1};
\node[idxlbl, anchor=east] at (-0.6,-1.800) {深度 2};
\node[idxlbl, anchor=east] at (-0.6,-2.700) {深度 3};
\node[idxlbl, anchor=east] at (-0.6,-3.600) {深度 4};
\end{tikzpicture}
\diagcaption{七個串的最佳合併樹。 方框是初始串，圓圈是一次合併，數字是合併出來的串
長。 總成本是圓圈的數字和 $13 + 25 + 43 + 55 + 83 + 138 = 357$，也等於
$\sum \ell_i d_i = 5 \cdot 4 + 8 \cdot 4 + 12 \cdot 3 + 18 \cdot 3 + 25 \cdot 3 +
30 \cdot 2 + 40 \cdot 2 = 357$。}
\label{fig:gr-merge-tree}
\end{minipage}
\end{center}

和第~15~章的做法比較：一趟一趟地把相鄰的兩個串兩兩合併（第 1 趟 R1+R2、R3+R4、
R5+R6，R7 輪空），三趟的成本是 $98 + 138 + 138 = 374$。 那種做法讓每個串的深度差不多，
串一樣長時很好；長短不一時，應該讓長的串少搬幾次。

\begin{examnote}
最佳合併的計算題，先看題目怎麼定義一次合併的成本。 本書用「搬動的紀錄數」$a + b$；
有的題目用最壞的比較次數 $a + b - 1$，$r$ 個串要合併 $r - 1$ 次，總成本差一個常數
$r - 1$，最佳的合併順序不變，但數字要照題目算。 問的是區塊 I/O 時，每次合併讀和寫各
$(a + b)/B$ 個區塊。 用公式 $\sum \ell_i d_i$ 驗算時，$d_i$ 是邊數（深度）；如果用
「根在第 $1$ 層」的層數代入，每一項都要減 $1$，否則多算了 $\sum \ell_i$。
\end{examnote}

\textbf{$k$ 路合併。} 每次合併 $k$ 個串，合併樹的每個內部節點都有 $k$ 個孩子。 引
理~\ref{lem:gr-merge-cost} 的證明完全不用改：總成本仍是 $\sum \ell_i d_i$。 直覺上
應該每次合併最短的 $k$ 個，第~15~章也是這麼說的，但這裡有一個細節。 每次合併讓串的
個數少 $k - 1$，從 $r$ 個變成 $1$ 個，如果每次都恰好合併 $k$ 個，$r - 1$ 必須是
$k - 1$ 的倍數。 不是倍數的時候，一定有某一次合併的串少於 $k$ 個，而這一次放在哪裡會影
響成本。

最好的做法是讓這次「不滿」的合併發生在最底下，處理最短的那幾個串。 實作上先補
\[
d = \bigl(k - 1 - (r - 1) \bmod (k - 1)\bigr) \bmod (k - 1)
\]
個長度 $0$ 的\textbf{虛擬串}，讓 $r + d - 1$ 是 $k - 1$ 的倍數，再每次合併最短的 $k$
個。 虛擬串長度是 $0$，搬它不花成本，它們一定在第一次合併就被取出，效果就是第一次只合
併最短的 $k - d$ 個真正的串。 這個做法的正確性和霍夫曼演算法的證明一樣，只是引
理~\ref{lem:gr-huff-choice} 改成「最小的 $k$ 個是同一個節點的孩子，在最深的一層」；
補虛擬串正是為了讓最深的那一層剛好有 $k$ 個兄弟。

補完之後串的個數 $r + d \equiv 1 \pmod{k - 1}$，每次合併少 $k - 1$ 個，所以每次都取得到
$k$ 個，最後剛好剩一個。 $k = 2$ 時 $d = 0$，就是霍夫曼演算法。

\textbf{例子。} 拿掉上面的 R7，剩下六個串 $12, 30, 8, 25, 5, 18$，用 $3$ 路合併。
$r = 6$，$(r - 1) \bmod 2 = 1$，補 $d = 1$ 個虛擬串。

\begin{center}
\begin{tabular}{lll}
\toprule
 & 補虛擬串 & 不補，直接每次取最短的三個 \\
\midrule
第 1 次 & $0 + 5 + 8 = 13$ & $5 + 8 + 12 = 25$ \\
第 2 次 & $12 + 13 + 18 = 43$ & $18 + 25 + 25 = 68$ \\
第 3 次 & $25 + 30 + 43 = 98$ & $30 + 68 = 98$（只剩兩個） \\
\midrule
總成本 & $154$ & $191$ \\
\bottomrule
\end{tabular}
\end{center}

不補虛擬串時，前兩次都合併了三個，最後只剩兩個串可以合併，「不滿」的那一次落在根。 和
補虛擬串的樹比較，R4 和 R1 的深度各多 $1$，多搬了 $25 + 12 = 37$ 千筆，正好是
$191 - 154$。 補了虛擬串，「不滿」的那一次在最底下，只牽涉最短的 $5$ 和 $8$。 補虛擬串的樹裡，R4、R2 深度 $1$，R1、R6 深度 $2$，R5、R3
深度 $3$，$\sum \ell_i d_i = 55 + 60 + 39 = 154$。

同樣的道理用在 $k$ 元霍夫曼碼（用 $0, 1, \ldots, k - 1$ 這 $k$ 個符號編碼）：先補頻率
$0$ 的虛擬符號，再每次合併頻率最小的 $k$ 個。

\section{排程}
\label{sec:gr-sched}

\subsection{最小化平均完成時間}
\label{sec:gr-spt}

實驗室只有一台 GPU 伺服器，一次只能跑一個訓練工作，一個工作開始了就要跑完。 早上
$9$ 點五位同學同時送來工作，各要跑的小時數是：

\begin{center}
\begin{tabular}{c*{5}{c}}
\toprule
工作 & A & B & C & D & E \\
\midrule
執行時間 $p_i$（小時） & 7 & 2 & 5 & 1 & 4 \\
\bottomrule
\end{tabular}
\end{center}

工作 $i$ 的\textbf{完成時間} $C_i$ 是它跑完的時刻（從 $9$ 點起算）。 每位同學都在等
自己的結果，我們想讓平均完成時間 $\frac{1}{n}\sum C_i$ 最小，也就是讓 $\sum C_i$ 最
小。 第~6~章的印表機佇列最後留了一個問題：先來先服務不一定讓平均等待最短，這一小節回
答它。

照送來的順序 A、B、C、D、E 跑：完成時間 $7, 9, 14, 15, 19$，總和 $64$，平均
$12.8$ 小時。 把短的排前面，D、B、E、C、A：完成時間 $1, 3, 7, 12, 19$，總和 $42$，平
均 $8.4$ 小時。 最後一個工作不管怎麼排都在 $19$ 點完成，差別在前面的工作：排在第 $k$
個的工作，它的執行時間會加進它自己和後面每一個工作的完成時間裡，被算 $n - k + 1$ 次。
寫成式子，順序是 $\pi(1), \pi(2), \ldots, \pi(n)$ 時
\[
\sum_{i=1}^{n} C_i = \sum_{k=1}^{n} (n - k + 1)\, p_{\pi(k)}.
\]
D、B、E、C、A 的順序代進去是 $5 \cdot 1 + 4 \cdot 2 + 3 \cdot 4 + 2 \cdot 5 + 1 \cdot 7 =
42$。 被算最多次的位置放最短的工作，這個貪心叫\textbf{最短工作優先}（shortest
processing time first，SPT）。

\begin{theorem}
\label{thm:gr-spt}
依執行時間由小到大排的順序，$\sum C_i$ 最小。 執行時間一樣的工作之間怎麼排，總和都
一樣。
\end{theorem}

\begin{proof}
任取一個順序。 如果它不是由小到大，一定有\textbf{相鄰}的兩個工作 $i$、$j$，$i$ 排在
$j$ 前面，$p_i > p_j$。 設它們前面的工作在時刻 $t$ 做完。 交換 $i$ 和 $j$ 之前，
$C_i = t + p_i$，$C_j = t + p_i + p_j$；交換之後，$C_j' = t + p_j$，$C_i' = t + p_j + p_i$。
其他工作的完成時間都不變（它們前面的工作總時數沒變）。 所以
\[
(C_i + C_j) - (C_i' + C_j') = p_i - p_j > 0,
\]
交換讓總和嚴格變小。 因此最佳的順序不能有這種相鄰的一對，也就是一定由小到大排好。
執行時間相同的工作交換，總和不變，所以所有由小到大的順序總和都一樣，都是最佳的。
\end{proof}

這個交換論證的寫法和活動選擇不太一樣：它說明任何不是由小到大的順序都還能再改進，沒有
去改造某一個最佳解。 兩種寫法都常見。

\textbf{等待時間。} 工作 $i$ 的等待時間是 $W_i = C_i - p_i$（開始之前等了多久），
$\sum W_i = \sum C_i - \sum p_i$，後面一項和順序無關，所以讓總等待時間最小的也是 SPT。
章首例子的 SPT 順序，總等待時間是 $42 - 19 = 23$ 小時。 演算法就是排序，
$\Theta(n \lg n)$。

工作有輕重之分（例如快要口試的同學比較急）時，每個工作有權重 $w_i$，要讓 $\sum w_i
C_i$ 最小。 同樣的相鄰交換論證可以算出應該依 $p_i / w_i$ 由小到大排，見觀念題
第~\ref{q:gr-wspt}~題。

\subsection{有期限的單位時間工作}
\label{sec:gr-deadline}

系上的雷射切割機開放預約，每個工作剛好用 $1$ 小時。 今天收到 $n$ 件申請，第 $i$ 件
的\textbf{期限}是 $d_i$（正整數）：要在第 $d_i$ 小時結束前做完，申請人才付 $p_i$ 元
（$p_i > 0$）；趕不上就不收費，乾脆不做。 時段編號 $1, 2, 3, \ldots$，時段 $t$ 是第
$t$ 個小時。 工作放在時段 $t$ 能趕上期限，條件是 $t \le d_i$。 要挑一些工作、各放進
不同的時段，讓收到的錢最多。

\begin{center}
\begin{tabular}{c*{9}{c}}
\toprule
工作 & 1 & 2 & 3 & 4 & 5 & 6 & 7 & 8 & 9 \\
\midrule
期限 $d_i$ & 3 & 1 & 4 & 5 & 4 & 1 & 6 & 3 & 2 \\
收費 $p_i$ & 45 & 20 & 60 & 35 & 50 & 25 & 15 & 40 & 30 \\
\bottomrule
\end{tabular}
\end{center}

\textbf{貪心的做法。} 依收費由高到低看每件工作（一樣時編號小的先），如果還找得到一個
不超過期限的空時段，就接下來，放進\textbf{最晚}的那個空時段；找不到就不接。 放最晚的
空時段，是為了把前面的時段留給期限比較早的工作。

期限超過 $n$ 的工作和期限是 $n$ 的沒有差別（$n$ 件工作用不到第 $n$ 個時段以後），可以
先把期限截到 $n$。 最直接的寫法是從期限往回一格一格找空時段，最壞每件工作走 $n$ 步，
總共 $O(n^2)$；第~\ref{sec:gr-dsu}~小節用互斥集合把它降到幾乎線性。

依收費排好是 $3, 5, 1, 8, 4, 9, 6, 2, 7$。 照這個做法，時段 $1$ 到 $6$ 依序放工作
$8, 1, 5, 3, 4, 7$，工作 $9, 6, 2$ 找不到位置，收費 $245$ 元；過程見下一小節的追蹤表。

\subsection{為什麼這樣排最多}

先把「一群工作排得進去」說清楚。 一群工作 $A$ 叫做\textbf{可行}的，是指能把它們放進
不同的時段，每件都不超過自己的期限。 令 $N_t(A)$ 是 $A$ 裡期限 $\le t$ 的工作件數。

\begin{lemma}
\label{lem:gr-feasible}
$A$ 是可行的，若且唯若對每個 $t \ge 1$，$N_t(A) \le t$。
\end{lemma}

\begin{proof}
必要：期限 $\le t$ 的工作都要放在時段 $1, \ldots, t$，而且各佔一個，所以最多 $t$
件。 充分：把 $A$ 依期限由小到大排，第 $k$ 件放時段 $k$。 如果第 $k$ 件的期限
$d < k$，那麼前 $k$ 件的期限都 $\le d$，$N_d(A) \ge k > d$，和條件矛盾。 所以每件都
趕得上。
\end{proof}

\begin{whybox}[貪心是最佳的]
證明分兩部分。

\textbf{(1) 找得到空時段，若且唯若接下來仍然可行。} 設目前接下的工作集合是 $A$（都已
經放好），輪到工作 $j$。 找得到空時段的話，直接得到 $A \cup \{j\}$ 的排法。 反過來，
假設時段 $1, \ldots, d_j$ 全滿了。 令 $u$ 是比 $d_j$ 大的第一個空時段（時段 $n + 1$
以後永遠是空的，所以 $u$ 存在；而且 $u \le n$，否則時段 $1, \ldots, n$ 全滿，$n$ 件工作都已
接下，不會輪到 $j$）。 時段 $1, \ldots, u - 1$ 全滿。 說明這些時段裡每件
工作的期限都 $\le u - 1$：假設某件工作 $q$ 放在時段 $s < u$，期限 $d_q \ge u$。 時段
不會被清空，所以 $q$ 被放進去的時候，時段 $u$ 是空的，而且 $u \le d_q$；演算法放的是
最晚的空時段，應該放在 $u$ 或更晚，不會是 $s$，矛盾。 所以 $A \cup \{j\}$ 裡期限
$\le u - 1$ 的工作至少有 $(u - 1) + 1$ 件（$d_j \le u - 1$），由
引理~\ref{lem:gr-feasible}，$A \cup \{j\}$ 不可行。

\textbf{(2) 交換論證。} 證明下面的不變量：每看完一件工作，已接下的集合 $A$ 都包含於
某個最佳解 $O$（收費最多的可行集合）。 一開始 $A$ 是空的，成立。 輪到工作 $j$ 時：

\begin{itemize}
\item 不接 $j$：由 (1)，$A \cup \{j\}$ 不可行；$O \supseteq A$，所以 $O$ 也不含 $j$
（可行集合的子集合還是可行的），不變量仍然成立。
\item 接下 $j$，而且 $j \in O$：不變量成立。
\item 接下 $j$，但 $j \notin O$：先說明 $O \setminus A$ 裡的工作都還沒輪到。 如果某件
$q \in O \setminus A$ 已經輪過、被拒絕了，那時接下的集合 $A_q \subseteq A \subseteq O$，
$A_q \cup \{q\} \subseteq O$ 是可行的，由 (1) 不會被拒絕，矛盾。 所以它們的收費都
$\le p_j$。 接著，$O \cup \{j\}$ 收費比 $O$ 多，不可能可行；令 $t^*$ 是讓
$N_t(O \cup \{j\}) > t$ 的最小 $t$。 $O$ 可行而多了一件，所以
$N_{t^*}(O \cup \{j\}) = t^* + 1$；$A \cup \{j\}$ 可行，所以 $N_{t^*}(A \cup \{j\})
\le t^*$。 兩者相差的工作都在 $O \setminus A$ 裡，因此有某件 $k \in O \setminus A$，
$d_k \le t^*$。 令 $O' = O \cup \{j\} \setminus \{k\}$。 對 $t < t^*$，
$N_t(O') \le N_t(O \cup \{j\}) \le t$；對 $t \ge t^*$，$k$ 被算在 $N_t$ 裡又被拿掉，
$N_t(O') = N_t(O \cup \{j\}) - 1 \le N_t(O) \le t$。 所以 $O'$ 可行，收費
$\ge$ $O$ 的收費（$p_k \le p_j$），也是最佳解，而且包含 $A \cup \{j\}$。
\end{itemize}

全部看完時，$A \subseteq O$ 而且 $O \setminus A$ 裡沒有工作（都輪過了，由上面的論證
不會被拒絕），所以 $A = O$，貪心的結果是最佳解。
\end{whybox}

\subsection{用互斥集合找最晚的空時段}
\label{sec:gr-dsu}

一格一格往回找太慢。 把時段分組：每一組
是一段連續的時段 $\{e, e + 1, \ldots, e + m\}$，$e$ 是空的（或是 $0$），後面 $m$ 個都
滿了。 時段 $0$ 是哨兵，代表「沒有空時段了」，它永遠不會被填滿。 這樣分組之後，「不
超過 $t$ 的最晚空時段」就是 $t$ 所在那一組的 $e$。 這正是第~13~章的互斥集合：

\begin{itemize}
\item 一開始每個時段 $0, 1, \ldots, D$ 自己一組（\proc{Make-Set}），$D = \min(n,
\max_i d_i)$。
\item 查詢：\proc{Find-Set}$(t)$ 找到 $t$ 那一組的代表，代表上記著這一組的空時段。
\item 時段 $e$ 被填滿以後，它那一組要併進左邊那一組（含 $e - 1$ 的那組）：
\proc{Union}$(e, e - 1)$。 合併後的空時段是左邊那一組原本的空時段。
\end{itemize}

第~13~章的 \proc{Union} 依大小合併，留下來的根不一定是空時段，所以另外用一個陣列
$\mathit{latest}$，只在根上有意義：$\mathit{latest}[r]$ 是根 $r$ 那一組的空時段。

\begin{pseudo}[label=alg:gr-deadline-dsu]{有期限的工作排程，用互斥集合（工作已依 $p$ 由大到小排好）}
\Function{Schedule-Jobs}{$d, p, n$}
    \State $D \gets \min(n, \max_{1 \le i \le n} d[i])$
    \For{$t \gets 0$ \To $D$}
        \State \Call{Make-Set}{$t$}
        \State $\mathit{latest}[t] \gets t$
        \State $\mathit{slot}[t] \gets \NIL$
    \EndFor
    \For{$i \gets 1$ \To $n$}
        \State $e \gets \mathit{latest}[\Call{Find-Set}{\min(d[i], D)}]$
        \If{$e > 0$} \Comment{$e = 0$ 表示沒有空時段，不接}
            \State $\mathit{slot}[e] \gets i$
            \State $\ell \gets \mathit{latest}[\Call{Find-Set}{e - 1}]$ \Comment{左邊那一組的空時段}
            \State \Call{Union}{$e, e - 1$}
            \State $\mathit{latest}[\Call{Find-Set}{e}] \gets \ell$
        \EndIf
    \EndFor
    \State \Return $\mathit{slot}$
\EndFunction
\end{pseudo}

第 11 行要在第 12 行的 \proc{Union} \textbf{之前}讀出左邊那一組的空時段；合併以後根可能換成右邊
那一組的根，就讀不到了。 每件工作做三次 \proc{Find-Set} 和一次 \proc{Union}，用依
大小合併加路徑壓縮，由第~13~章，總時間 $O(n\,\alpha(n))$，整個演算法被排序的
$O(n \lg n)$ 主導。

\textbf{追蹤。} 同樣的九件工作，$D = \min(9, 6) = 6$，時段 $0$ 到 $6$。 \proc{Union}
照第~13~章的規則：依大小合併，一樣大時第二個參數的根接到第一個參數的根底下；
\proc{Find-Set} 做路徑壓縮。 「$p[0..6]$」是做完這件工作之後的父節點陣列，
「$\mathit{latest}$」只列出各個根上的值（根:空時段）。

\begin{center}
\small
\begin{tabular}{cccccl}
\toprule
工作 & $\min(d, D)$ & 根 & 空時段 $e$ & $p[0..6]$ & $\mathit{latest}$（根:空時段） \\
\midrule
開始 & & & & 0 1 2 3 4 5 6 & 每個時段自己 \\
3 & 4 & 4 & 4 & 0 1 2 4 4 5 6 & 0:0, 1:1, 2:2, 4:3, 5:5, 6:6 \\
5 & 4 & 4 & 3 & 0 1 4 4 4 5 6 & 0:0, 1:1, 4:2, 5:5, 6:6 \\
1 & 3 & 4 & 2 & 0 4 4 4 4 5 6 & 0:0, 4:1, 5:5, 6:6 \\
8 & 3 & 4 & 1 & 4 4 4 4 4 5 6 & 4:0, 5:5, 6:6 \\
4 & 5 & 5 & 5 & 4 4 4 4 4 4 6 & 4:0, 6:6 \\
9 & 2 & 4 & 0，不接 & 不變 & \\
6 & 1 & 4 & 0，不接 & 不變 & \\
2 & 1 & 4 & 0，不接 & 不變 & \\
7 & 6 & 6 & 6 & 4 4 4 4 4 4 4 & 4:0 \\
\bottomrule
\end{tabular}
\end{center}

第一件工作 3 填了時段 $4$，\proc{Union}$(4, 3)$ 兩組一樣大，$3$ 接到 $4$ 底下，根是
$4$；但這一組的空時段是 $3$，所以 $\mathit{latest}[4] = 3$。 這就是為什麼不能直接拿根
當空時段。 工作 4 填了時段 $5$ 之後，\proc{Union}$(5, 4)$：$5$ 那一組只有 $1$ 個時段，
$4$ 那一組有 $5$ 個，依大小合併讓 $5$ 接到 $4$ 底下，$\mathit{latest}[4]$ 仍是 $0$。
從此期限 $\le 5$ 的工作一律一次 \proc{Find-Set} 就知道沒有空位。

另一種常見的寫法不依大小合併，填滿時段 $e$ 時直接設 $p[e] \gets \proc{Find-Set}(e - 1)$，
根永遠是空時段，不需要 $\mathit{latest}$；只靠路徑壓縮，總時間 $O(n \lg n)$。

\subsection{C 程式}

\begin{cbookrecap}[第~14~章結構、第~19~章 19.3 節]
\texttt{calloc(count, size)} 和 \texttt{malloc} 一樣要記憶體，但會把整塊設成 $0$；
下面的 \texttt{slot} 用 $0$ 表示「這個時段沒有工作」，工作從 $1$ 編號，所以剛好。 互
斥集合的 \texttt{dsu\_init}、\texttt{find\_iter}、\texttt{dsu\_union} 沿用第~13~章
範例 dsu\_trace.c 的寫法，只拿掉這裡用不到的欄位。
\end{cbookrecap}

程式讀進 $n$ 件工作，每件一列：期限和收費。 每接下一件工作，印出它放的時段和這時的
父節點陣列；最後印出時段 $1$ 到 $D$ 各放了哪件工作（\texttt{-} 是空的）和總收費。

\begin{example}[label=ex:gr-deadline]{c}{(deadline.c) 有期限的工作排程：用互斥集合找最晚的空時段}
#include <stdio.h>
#include <stdlib.h>

// ---- 第 13 章的互斥集合：依大小合併、迴圈版路徑壓縮（只留用得到的欄位）----
typedef struct {
    int *parent;    // parent[x] == x 表示 x 是根
    int *size;      // 只有根的 size 有意義
} DSU;

int dsu_init(DSU *d, int n)
{
    d->parent = malloc(n * sizeof *d->parent);
    d->size = malloc(n * sizeof *d->size);
    if (d->parent == NULL || d->size == NULL) {
        return 0;
    }
    for (int x = 0; x < n; x++) {
        d->parent[x] = x;
        d->size[x] = 1;
    }
    return 1;
}

int find_iter(int *parent, int x)
{
    int root = x;
    while (parent[root] != root) {
        root = parent[root];
    }
    while (parent[x] != root) {
        int next = parent[x];
        parent[x] = root;
        x = next;
    }
    return root;
}

void dsu_union(DSU *d, int x, int y)    // 一樣大時 y 的根接到 x 的根底下
{
    int rx = find_iter(d->parent, x), ry = find_iter(d->parent, y);
    if (rx == ry) {
        return;
    }
    if (d->size[rx] < d->size[ry]) {
        int t = rx;
        rx = ry;
        ry = t;
    }
    d->parent[ry] = rx;
    d->size[rx] += d->size[ry];
}

// ---- 有期限的單位時間工作 ----
typedef struct {
    int id, d, p;   // 編號、期限、收費
} Job;

int by_profit(const void *a, const void *b)    // 收費高的先；一樣時編號小的先
{
    const Job *x = a, *y = b;
    if (x->p != y->p) {
        return x->p > y->p ? -1 : 1;
    }
    return (x->id > y->id) - (x->id < y->id);
}

int main(void)
{
    int n, D = 0;
    if (scanf("%d", &n) != 1 || n <= 0) {
        return 1;
    }
    Job *job = malloc(n * sizeof *job);
    int ok = job != NULL;
    for (int i = 0; ok && i < n; i++) {
        ok = scanf("%d %d", &job[i].d, &job[i].p) == 2 && job[i].d >= 1;
        job[i].id = i + 1;
        if (ok && job[i].d > D) {
            D = job[i].d;
        }
    }
    D = D < n ? D : n;      // 時段 1..D；時段 0 是「沒有空位」的哨兵
    DSU ds = {NULL, NULL};
    int *latest = malloc((D + 1) * sizeof *latest);  // latest[根]：這組的空時段
    int *slot = calloc(D + 1, sizeof *slot);         // slot[t]：放在時段 t 的工作
    ok = ok && latest != NULL && slot != NULL && dsu_init(&ds, D + 1);
    if (ok) {
        for (int t = 0; t <= D; t++) {
            latest[t] = t;
        }
        qsort(job, n, sizeof job[0], by_profit);
        long long total = 0;
        for (int i = 0; i < n; i++) {
            int t = job[i].d < D ? job[i].d : D;
            int e = latest[find_iter(ds.parent, t)];
            printf("job %d (d=%d, p=%d): ", job[i].id, job[i].d, job[i].p);
            if (e == 0) {
                printf("rejected\n");
                continue;
            }
            slot[e] = job[i].id;
            total += job[i].p;
            int left = latest[find_iter(ds.parent, e - 1)];  // 合併前先讀
            dsu_union(&ds, e, e - 1);
            latest[find_iter(ds.parent, e)] = left;
            printf("slot %d, parent =", e);
            for (int x = 0; x <= D; x++) {
                printf(" %d", ds.parent[x]);
            }
            printf("\n");
        }
        printf("schedule:");
        for (int t = 1; t <= D; t++) {
            printf(slot[t] ? " %d" : " -", slot[t]);
        }
        printf("\ntotal profit = %lld\n", total);
    }
    free(job);
    free(latest);
    free(slot);
    free(ds.parent);
    free(ds.size);
    return ok ? 0 : 1;
}
\end{example}
\iosample{9 \\ 3 45 \\ 1 20 \\ 4 60 \\ 5 35 \\ 4 50 \\ 1 25 \\ 6 15 \\ 3 40 \\ 2 30}{job 3 (d=4, p=60): slot 4, parent = 0 1 2 4 4 5 6 \\ job 5 (d=4, p=50): slot 3, parent = 0 1 4 4 4 5 6 \\ job 1 (d=3, p=45): slot 2, parent = 0 4 4 4 4 5 6 \\ job 8 (d=3, p=40): slot 1, parent = 4 4 4 4 4 5 6 \\ job 4 (d=5, p=35): slot 5, parent = 4 4 4 4 4 4 6 \\ job 9 (d=2, p=30): rejected \\ job 6 (d=1, p=25): rejected \\ job 2 (d=1, p=20): rejected \\ job 7 (d=6, p=15): slot 6, parent = 4 4 4 4 4 4 4 \\ schedule: 8 1 5 3 4 7 \\ total profit = 245}

每一列的父節點陣列都和手算的表一樣。 程式的元素直接用時段編號 $0..D$，不必像
第~13~章那樣在輸入輸出時加減一。 \texttt{slot} 和 \texttt{latest} 都要 $D + 1$ 格，
因為時段 $0$ 也是一個元素。

\section{貪心何時不對}
\label{sec:gr-fail}

前面每一個貪心都附了證明。 這一節反過來，看幾個「看起來很合理、其實是錯的」貪心，最
後介紹把貪心整理成一般理論的擬陣。

\subsection{找零錢}

收銀機有幾種面額的硬幣 $c_1 < c_2 < \cdots < c_m$，每種都用不完，要湊出金額 $x$，硬
幣枚數越少越好。 直覺的貪心是：能用大的就用大的。

\begin{pseudo}[label=alg:gr-change]{找零錢的貪心}
\Function{Greedy-Change}{$c, m, x$}
    \State $\mathit{count} \gets 0$
    \For{$i \gets m$ \DownTo $1$}
        \State $\mathit{count} \gets \mathit{count} + \lfloor x / c[i] \rfloor$
        \State $x \gets x \bmod c[i]$
    \EndFor
    \If{$x > 0$}
        \State \Return \NIL \Comment{湊不出來}
    \EndIf
    \State \Return $\mathit{count}$
\EndFunction
\end{pseudo}

新臺幣常用的硬幣 $1, 5, 10, 50$ 元用這個貪心，每個金額都得到最少枚數（觀念題
第~\ref{q:gr-coin}~題會說明怎麼驗證）。 這樣的面額組合叫\textbf{標準的}（canonical）。
換一組面額就不一定了：

\begin{itemize}
\item 面額 $1, 5, 12$，金額 $15$：貪心先拿 $12$，剩 $3$ 只能用三個 $1$，共 $4$ 枚；
三個 $5$ 只要 $3$ 枚。
\item 面額 $4, 10$（沒有 $1$），金額 $12$：貪心先拿 $10$，剩 $2$ 湊不出來，回傳
\NIL；其實 $4 + 4 + 4$ 就可以。 貪心不只不是最少，連有解都找不到。
\end{itemize}

出錯的原因和 0/1 背包一樣：拿了 $12$ 之後剩下的 $3$ 很難湊，貪心在選 $12$ 的時候看不
到這件事。 一般的面額要用動態規劃：令 $M[a]$ 是湊出 $a$ 的最少枚數，$M[0] = 0$，
$M[a] = 1 + \min_{c_i \le a} M[a - c_i]$，時間 $O(mx)$，是擬多項式的（第~18~章）。

一組面額是不是標準的，也有辦法判斷。 Kozen 與 Zaks（1994）證明：如果貪心對某個金額
不是最少枚數，那麼最小的反例一定小於 $c_{m-1} + c_m$。 所以只要用動態規劃算出這個範
圍內每個金額的最少枚數，逐一和貪心比較，就能確定這組面額是不是標準的。

\subsection{其他常見的錯誤貪心}

\begin{itemize}
\item \textbf{0/1 背包}：第~\ref{sec:gr-knapsack}~節。 不管依單位重量價值、價值還是
重量，都有反例。
\item \textbf{加權的活動選擇}：每筆申請有不同的租金，要讓收入最多。 最早結束的先排會
錯：$[8, 10)$ 租金 $1$、$[9, 12)$ 租金 $10$，它先選了前者，後者就排不進去。 這時要用
動態規劃：依結束時刻排好，第 $i$ 筆「選或不選」，選的話接上最後一筆和它相容的申請。
\item \textbf{期限排程放最早的空時段}：第~\ref{sec:gr-deadline}~節如果把工作放在最早
的空時段：工作 A（期限 $2$、收費 $10$）先放進時段 $1$，工作 B（期限 $1$、收費 $8$）就
沒位置了，只收 $10$；放最晚的空時段可以兩件都接，收 $18$。 同一個演算
法，細節換一個，就不對了。
\item 第~25~章的 Dijkstra 演算法也是貪心，邊的權重有負的時候就會錯。
\end{itemize}

遇到一個新的貪心想法，先用小例子找反例，特別是讓兩個選擇平手或差一點點的例子；找不
到，再試著寫交換論證，寫不下去的地方往往就是反例的方向。 也可以另外寫一個窮舉的版本，
在大量隨機的小輸入上和貪心比對答案。 比對一萬組都相同也不是證明，但只要有一組不同，貪
心就是錯的。

\subsection{擬陣：貪心的一般理論（觀念）}
\label{sec:gr-matroid}

很多貪心都長成同一個樣子：有一群元素，每個有正的權重；有一些子集合叫做「獨立的」
（可以同時選）；依權重由大到小看每個元素，加進去之後仍然獨立就加。 第~\ref{sec:gr-deadline}~節的期限排程正是這樣：元素是工作，權重是收費，「獨立」是「可行」。 這種貪心什麼時候對？
有一個漂亮的答案。

一個有限集合 $S$ 和一群「獨立」的子集合 $\mathcal{I}$，如果滿足下面三條，叫做
\textbf{擬陣}（matroid）：

\begin{enumerate}
\item 空集合是獨立的。
\item \textbf{遺傳性}：獨立集合的子集合也是獨立的。
\item \textbf{交換性}：$A$、$B$ 都是獨立的，$|A| < |B|$，那麼 $B$ 裡一定有某個不在
$A$ 裡的元素 $x$，讓 $A \cup \{x\}$ 仍然獨立。
\end{enumerate}

\textbf{定理}（Rado 1957、Edmonds 1971）：在擬陣上，對任何正的權重，上面的貪心都得到
權重和最大的獨立集合。 反過來，如果 $(S, \mathcal{I})$ 滿足前兩條，而且對每一種權重
貪心都對，它一定也滿足第三條，是擬陣。 這裡不證明。

幾個例子：

\begin{itemize}
\item \textbf{圖的擬陣}：元素是無向圖的邊，「獨立」是「不含環」。 權重由小到大的版本就是
第~24~章的 Kruskal 最小生成樹演算法。
\item \textbf{期限排程}：元素是工作，「獨立」是「可行」；交換性可以用
引理~\ref{lem:gr-feasible} 證明。 第~\ref{sec:gr-deadline}~節的演算法就是擬陣上的貪心，「找
得到空時段」就是檢查加進去之後仍然獨立。
\item \textbf{不是擬陣的例子}：兩兩相容的區間。 $A = \{[1, 10)\}$、$B = \{[1, 2), [3, 4)\}$，
$B$ 的兩個都和 $[1, 10)$ 衝突，交換性不成立。 活動選擇的貪心卻是對的，因為它依的是結束
時刻這個特別的順序，不是權重。 擬陣是貪心成立的充分條件，不是必要條件；霍夫曼編碼也不
是擬陣上的問題。
\end{itemize}

CLRS 第 4 版把擬陣的完整內容移到了書的網站，第 3 版的 16.4、16.5 節有證明和用擬陣
處理期限排程的例子；第 4 版第 15 章的章末附注列了原始文獻。

\section{本章常見錯誤}
\label{sec:gr-pitfalls}

\begin{center}
\begin{tabular}{p{0.24\linewidth}p{0.50\linewidth}p{0.16\linewidth}}
\toprule
錯誤 & 例子 & 見 \\
\midrule
排序的鍵選錯 & 活動選擇依開始時刻或長度排序；以為「衝突最少的先排」一定對 &
第~\ref{sec:gr-activity}~節 \\
端點的相容條件弄錯 & 半開區間寫成 $s_i > \mathit{last}$，少選了首尾相接的申請；閉區
間反而寫成 $\ge$ & 第~\ref{sec:gr-activity}~節 \\
交換論證只換一次 & 證明了第一個選擇對，就說整個貪心對，沒有用最佳子結構或歸納接下去；
或換完沒檢查解還可行 & 第~\ref{sec:gr-elements}~節 \\
分數背包的做法用在 0/1 背包 & 依單位重量價值整件拿，袋子剩下的空間浪費掉 &
第~\ref{sec:gr-knapsack}~節 \\
碼長多算 $1$ & 把碼長寫成葉節點所在的層數；碼長是深度（邊數） &
第~\ref{sec:gr-huffman}~節 \\
$k$ 路合併沒補虛擬串 & $r - 1$ 不是 $k - 1$ 的倍數，直接每次取最短的 $k$ 個 &
第~\ref{sec:gr-merge}~節 \\
放到最早的空時段 & 期限排程把工作放進最早的空時段，擋住後面期限早的工作 &
第~\ref{sec:gr-deadline}~節 \\
把根當成空時段 & 互斥集合用了依大小合併，卻直接拿 \proc{Find-Set} 的結果當空時段；
或合併之後才去讀左邊那組的空時段 & 第~\ref{sec:gr-deadline}~節 \\
\bottomrule
\end{tabular}
\end{center}

\section{觀念題}
\label{sec:gr-exercises}

題型和研究所考試、期中期末考的筆試相近，有選擇、計算、追蹤和證明。 解答在本章最後。

\begin{enumerate}
\item \label{q:gr-act-trace}（追蹤）籃球場收到十筆借用申請，時段都是半開區間：
\begin{center}
\begin{tabular}{c*{10}{c}}
\toprule
編號 & 1 & 2 & 3 & 4 & 5 & 6 & 7 & 8 & 9 & 10 \\
\midrule
$s_i$ & 6 & 1 & 2 & 4 & 8 & 5 & 9 & 11 & 3 & 10 \\
$f_i$ & 9 & 3 & 6 & 6 & 10 & 8 & 12 & 13 & 4 & 12 \\
\bottomrule
\end{tabular}
\end{center}
\begin{enumerate}
\item 用虛擬碼~\ref{alg:gr-activity}（結束時刻相同時編號小的先）寫出排序後的順序和選中
的申請。
\item 結束時刻相同時改成編號大的先，選中的申請有沒有變？ 筆數呢？
\item 如果時段改成閉區間（首尾相接也算衝突），最多能選幾筆？
\end{enumerate}

\item （選擇，可能不只一個答案）下列哪些敘述正確？
\begin{enumerate}
\item 分數背包依單位重量價值由大到小拿，得到最大總價值。
\item 0/1 背包依單位重量價值由大到小拿（放得下才拿），得到的總價值至少是最佳值的一半。
\item 霍夫曼演算法建出的碼樹一定是嚴格二元樹。
\item 頻率嚴格最高（比其他每個符號都高）的符號，在任何最佳前綴碼裡，碼長都不比其他
符號長。
\item $2$ 路合併時，每次合併目前最短的兩個串，總合併成本最小。
\item $3$ 路合併時，每次合併目前最短的三個串，總合併成本一定最小。
\item 一個最佳化問題只要有最佳子結構，貪心就一定對。
\end{enumerate}

\item \label{q:gr-huff-trace}（追蹤）八個符號的頻率如下。 用本書的平手規則（先比頻率、再比編號；葉節點照表上
的順序編號 $0$ 到 $7$；先取出的當左孩子，左邊標 \texttt{0}）建霍夫曼樹，寫出每一次合
併、每個符號的碼和總位元數，並和定長碼比較。
\begin{center}
\begin{tabular}{c*{8}{c}}
\toprule
符號 & a & b & c & d & e & f & g & h \\
\midrule
頻率 & 9 & 4 & 15 & 4 & 30 & 7 & 12 & 19 \\
\bottomrule
\end{tabular}
\end{center}

\item （選擇，可能不只一個答案）六個符號，頻率都是正的。 下列哪些碼長組合，可能是某
組頻率的最佳前綴碼的碼長？
\begin{enumerate}
\item $1, 2, 3, 4, 5, 5$
\item $2, 2, 2, 3, 3, 3$
\item $2, 2, 3, 3, 3, 3$
\item $2, 2, 2, 3, 4, 4$
\item $1, 3, 3, 3, 3, 3$
\item $2, 2, 3, 3, 3, 4$
\end{enumerate}

\item \label{q:gr-kmerge-q}（計算）外部排序產生了八個串，長度（千筆）依序是
$14, 6, 22, 9, 30, 11, 3, 17$。 一次合併的成本是合併出來的串長。
\begin{enumerate}
\item 求 $2$ 路合併的最小總成本，寫出每一次合併。
\item 改成 $3$ 路合併：要補幾個虛擬串？ 最小總成本是多少？
\item $3$ 路合併但不補虛擬串，每次直接合併最短的三個（最後不夠三個就有幾個合幾個），
總成本是多少？
\end{enumerate}

\item \label{q:gr-wspt}（證明）同第~\ref{sec:gr-spt}~節，但工作 $i$ 有權重 $w_i > 0$，
要讓 $\sum w_i C_i$ 最小。 證明：依 $p_i / w_i$ 由小到大排的順序是最佳的。 所有
$w_i = 1$ 時，這個結果變成什麼？

\item \label{q:gr-dl-trace}（追蹤）八件單位時間的工作：
\begin{center}
\begin{tabular}{c*{8}{c}}
\toprule
工作 & 1 & 2 & 3 & 4 & 5 & 6 & 7 & 8 \\
\midrule
期限 & 2 & 4 & 1 & 3 & 2 & 4 & 5 & 3 \\
收費 & 40 & 15 & 30 & 55 & 25 & 45 & 10 & 20 \\
\bottomrule
\end{tabular}
\end{center}
用虛擬碼~\ref{alg:gr-deadline-dsu}（\proc{Union} 依大小合併，一樣大時第二個參數的根接
到第一個的根底下；\proc{Find-Set} 做路徑壓縮）排程。 寫出每接下一件工作後的父節點陣
列 $p[0..D]$，以及最後的排程和總收費。

\item \label{q:gr-coin}（計算）
\begin{enumerate}
\item 面額 $1, 4, 9$。 找出貪心（虛擬碼~\ref{alg:gr-change}）不是最少枚數的最小金額，
寫出貪心的結果和最佳的湊法。
\item 用 Kozen 與 Zaks 的結論，要確定 $1, 5, 10, 50$ 是標準的，需要檢查哪些金額？
說明這組面額為什麼是標準的。
\end{enumerate}

\item \label{q:gr-ahead}（證明）另一種證明活動選擇的方法。 設貪心選到的申請依序是
$g_1, g_2, \ldots, g_k$，$o_1, o_2, \ldots, o_m$ 是任何一組兩兩相容的申請，依結束時刻
排好。 證明：對每個 $i \le \min(k, m)$，$f(g_i) \le f(o_i)$。 再用這個結果證明
$k \ge m$。

\item （選擇，可能不只一個答案）下列哪些是擬陣（元素集合／獨立的子集合）？
\begin{enumerate}
\item 一個無向圖的邊／不含環的邊集合。
\item 一群區間／兩兩相容（不重疊）的區間集合。
\item 一群有期限的單位時間工作／能在期限內全部做完的工作集合。
\item 一個無向圖的邊／沒有兩條邊共用端點的邊集合（匹配）。
\item 一個集合 $S$ 的元素／元素個數不超過 $3$ 的子集合。
\end{enumerate}
\end{enumerate}

\section{實作練習}
\label{sec:gr-practice}

下面是 Aki-Judge 上和本章相關的題目。 題目不多，霍夫曼編碼、最佳合併和期限排程目前
沒有對應的實作題，請用觀念題第~\ref{q:gr-huff-trace}--\ref{q:gr-kmerge-q}、\ref{q:gr-dl-trace}~題和範例~\ref{ex:gr-huffman}、
範例~\ref{ex:gr-deadline} 的程式練習。

\begin{center}
\begin{tabular}{llc}
\toprule
小節 & 題目 & 難度 \\
\midrule
活動選擇 & \akiref{most-compatible-bookings}{最多筆不衝突的借用申請} & 基本 \\
分數背包 & \akiref{maximum-value-with-divisible-items}{可分割物品的最大總價值} & 基本 \\
最小化完成時間 & \akiref{smallest-total-waiting-time}{最小的總等待時間} & 暖身 \\
貪心何時不對 & \akiref{fewest-coins-for-an-amount}{湊出金額的最少硬幣數} & 進階 \\
\bottomrule
\end{tabular}
\end{center}

幾題的提示：

\begin{itemize}
\item 「最多筆不衝突的借用申請」：時段是半開區間，和範例~\ref{ex:gr-activity} 一樣；
$n$ 可以是 $0$。 只要印筆數。
\item 「可分割物品的最大總價值」：比值用交叉相乘比較，$v_i w_j$ 對 $v_j w_i$，不要用
浮點數除法；乘積要用 \texttt{long long}。 拿一部分時的價值是整數除法、無條件捨去，
題目也規定了比值相同時的先後。
\item 「最小的總等待時間」問的是等待時間（開始之前等了多久），不是完成時間；總和可能
超過 \texttt{int} 的範圍。
\item 「湊出金額的最少硬幣數」的面額是任意的，貪心會錯，要用第~\ref{sec:gr-fail}~節
的遞迴式做動態規劃（第~18~章）。 先用 $1, 5, 12$ 湊 $15$ 這類例子確認你的程式不是貪
心。
\end{itemize}

\section{讀完你應該會}
\label{sec:gr-checklist}

\begin{itemize}
\item 對一個貪心演算法，寫出「貪心選擇性質」和「最佳子結構」兩部分的證明；會用交換論
證，也會用逐步比較（觀念題第~\ref{q:gr-ahead}~題）的寫法。
\item 手算活動選擇、分數背包，並舉出 0/1 背包、找零錢、加權活動選擇讓貪心失敗的反例。
\item 依指定的平手規則逐步畫出霍夫曼樹，寫出每個符號的碼、總位元數，並解碼一段位元
串；知道碼長是葉節點的深度（邊數）。
\item 算出 $2$ 路和 $k$ 路的最佳合併成本，知道什麼時候要補幾個虛擬串，並用
$\sum \ell_i d_i$ 驗算。
\item 用最短工作優先排程並證明它最佳；用互斥集合實作有期限的工作排程，手算每一步的父
節點陣列。
\item 說出擬陣的定義，判斷一個「獨立集合」的系統是不是擬陣。
\end{itemize}

\textbf{延伸閱讀}（書末會整理完整的參考資料）

\begin{itemize}
\item Thomas H. Cormen 等，\textit{Introduction to Algorithms}，第 4 版，第 15 章：活
動選擇、貪心法的要素、分數與 0/1 背包、霍夫曼碼，證明比本章更形式化。 擬陣和用擬陣
處理期限排程在第 3 版的 16.4、16.5 節（第 4 版移到書的網站）。
\item Jon Kleinberg 與 Éva Tardos，\textit{Algorithm Design}，第 4 章：把貪心的證明分成
「逐步領先」和「交換論證」兩種寫法，例子多，適合想練習寫證明的人。
\item Steven S. Skiena，\textit{The Algorithm Design Manual}，第 3 版，第 1 章 1.2、
1.3 節：用排程問題示範怎麼找貪心的反例。
\item Antti Laaksonen，\textit{Competitive Programmer's Handbook}，第 6 章：找零、排
程、期限、霍夫曼碼，每個主題一兩頁，適合快速複習。
\item Donald E. Knuth，\textit{The Art of Computer Programming}，第 1 卷，第 3 版，
2.3.4.5 節：路徑長與霍夫曼的方法，包括 $t$ 元樹要補權重 $0$ 的節點。
\end{itemize}

\addcontentsline{toc}{section}{觀念題解答}
\section*{觀念題解答}

\begin{enumerate}
\item
\begin{enumerate}
\item 依結束時刻排好（平手編號小的先）：$2, 9, 3, 4, 6, 1, 5, 7, 10, 8$。 選 $2$
$[1, 3)$；$9$ $[3, 4)$ 在 $3$ 開始，選；$3$ $[2, 6)$ 跳過；$4$ $[4, 6)$ 選；$6$
$[5, 8)$ 跳過；$1$ $[6, 9)$ 選；$5$ $[8, 10)$ 跳過；$7$ $[9, 12)$ 選；$10$、$8$ 跳過。
選中 $2, 9, 4, 1, 7$，共 $5$ 筆。
\item 順序變成 $2, 9, 4, 3, 6, 1, 5, 10, 7, 8$。 $4$ 和 $3$ 對調不影響（$3$ 仍然被跳
過）；$10$ $[10, 12)$ 排到 $7$ 前面，它在 $10 \ge 9$ 開始，被選中，$7$ 改被跳過。 選中
$2, 9, 4, 1, 10$，換了一筆，筆數仍是 $5$（定理~\ref{thm:gr-act-choice} 對任何一筆結束最早
的都成立）。
\item 閉區間時相容的條件是 $s_i > \mathit{last}$。 順序同 (a)：選 $2$ $[1, 3]$；$9$
在 $3$ 開始，不行；$3$ 不行；$4$ $[4, 6]$ 選；$6$ 不行；$1$ 在 $6$ 開始，不行；$5$
$[8, 10]$ 選；$7$、$10$ 不行；$8$ $[11, 13]$ 選。 共 $4$ 筆；證明照樣成立，所以最多就是 $4$ 筆。
\end{enumerate}

\item (a)、(c)、(d)、(e) 正確。 (b)：容量 $10$，物品（重 $1$，值 $2$）和（重 $10$，
值 $10$）。 前者比值 $2$ 先拿，剩 $9$ 放不下後者，得 $2$，最佳是 $10$，連一半都不到；
把後者的重量和價值放大，比例可以任意小。 (c)：每次合併都接兩個孩子，沒有只有一個孩
子的節點。 (d)：如果 $x$ 頻率比 $y$ 高卻比 $y$ 深，交換兩者 $B$ 嚴格變小（第~\ref{sec:gr-huff-proof}~節的式子），和最佳矛盾。 (f)：第~\ref{sec:gr-merge}~節
的例子，不補虛擬串得到 $191$，最佳是 $154$。 (g)：0/1 背包有最佳子結構，貪心卻會錯。

\item 編號：a $0$、b $1$、c $2$、d $3$、e $4$、f $5$、g $6$、h $7$，內部節點從 $8$
開始。 七次合併（先取出的寫在前面）：b $4$ + d $4$ = \#8 $8$；f $7$ + \#8 $8$ = \#9 $15$；
a $9$ + g $12$ = \#10 $21$；c $15$ + \#9 $15$ = \#11 $30$（平手，葉節點編號小，先取 c）；
h $19$ + \#10 $21$ = \#12 $40$；e $30$ + \#11 $30$ = \#13 $60$（平手，先取 e）；\#12 + \#13 = $100$。
根的左子樹是 \#12（h 和 \#10），右子樹是 \#13（e 和 \#11）。 碼：h \texttt{00}、
a \texttt{010}、g \texttt{011}、e \texttt{10}、c \texttt{110}、f \texttt{1110}、
b \texttt{11110}、d \texttt{11111}。 總位元數
$9 \cdot 3 + 4 \cdot 5 + 15 \cdot 3 + 4 \cdot 5 + 30 \cdot 2 + 7 \cdot 4 + 12 \cdot 3
+ 19 \cdot 2 = 274$，也等於內部節點的和 $8 + 15 + 21 + 30 + 40 + 60 + 100 = 274$。
定長碼要 $3 \times 100 = 300$ 位元。 把範例~\ref{ex:gr-huffman} 的輸入改成這八個符號，
可以核對每一個碼。

\item (a)、(c)、(d)。 嚴格二元樹的葉節點深度一定滿足 $\sum 2^{-l_i} = 1$（根算 $1$，每個內
部節點把自己的值平分給兩個孩子），反過來滿足這個等式的碼長都排得出一棵嚴格二元樹。 (a)
$\frac12 + \frac14 + \frac18 + \frac1{16} + \frac2{32} = 1$；(b) $\frac34 + \frac38 > 1$，連
前綴碼都不是；(c) $\frac12 + \frac12 = 1$；(d) $\frac34 + \frac18 + \frac2{16} = 1$；(e)
$\frac12 + \frac58 > 1$；(f) $\frac{15}{16} < 1$，碼樹不是嚴格的，不是最佳。 (a)、(c)、(d)
分別是頻率 $16, 8, 4, 2, 1, 1$、全部相同、$4, 4, 4, 2, 1, 1$ 的霍夫曼碼長。

\item
\begin{enumerate}
\item 每次合併最短的兩個：$3 + 6 = 9$；$9 + 9 = 18$（原本的 $9$ 和剛合出的 $9$）；
$11 + 14 = 25$；$17 + 18 = 35$；$22 + 25 = 47$；$30 + 35 = 65$；$47 + 65 = 112$。
總成本 $311$。 （照第~15~章一趟一趟兩兩相鄰合併，三趟共 $336$。）
\item $r = 8$，$(8 - 1) \bmod 2 = 1$，補 $1$ 個虛擬串。 $0 + 3 + 6 = 9$；
$9 + 9 + 11 = 29$；$14 + 17 + 22 = 53$；$29 + 30 + 53 = 112$。 總成本 $203$。
\item $3 + 6 + 9 = 18$；$11 + 14 + 17 = 42$；$18 + 22 + 30 = 70$；最後只剩兩個，
$42 + 70 = 112$。 總成本 $242$，比 (b) 多 $39$。
\end{enumerate}

\item 設相鄰的 $i$、$j$，$i$ 在前，前面的工作在 $t$ 做完。 交換前
$w_i(t + p_i) + w_j(t + p_i + p_j)$，交換後 $w_j(t + p_j) + w_i(t + p_j + p_i)$，其他
工作不變。 相減：
\[
(\text{交換前}) - (\text{交換後}) = w_j p_i - w_i p_j .
\]
若 $p_i / w_i > p_j / w_j$，即 $p_i w_j > p_j w_i$，交換讓總和嚴格變小。 所以最佳順序
不會有相鄰的一對違反「$p / w$ 由小到大」，也就是一定照 $p / w$ 由小到大排；比值相同
的交換總和不變。 $w_i$ 全部是 $1$ 時，$p / w = p$，就是定理~\ref{thm:gr-spt} 的最短工
作優先。

\item $D = \min(8, 5) = 5$。 依收費排好：$4, 6, 1, 3, 5, 8, 2, 7$。
\begin{center}
\begin{tabular}{cccl}
\toprule
工作 & $\min(d, D)$ & 結果 & $p[0..5]$ \\
\midrule
4 & 3 & 時段 3 & 0 1 3 3 4 5 \\
6 & 4 & 時段 4 & 0 1 3 3 3 5 \\
1 & 2 & 時段 2 & 0 3 3 3 3 5 \\
3 & 1 & 時段 1 & 3 3 3 3 3 5 \\
5 & 2 & 不接 & 不變 \\
8 & 3 & 不接 & 不變 \\
2 & 4 & 不接 & 不變 \\
7 & 5 & 時段 5 & 3 3 3 3 3 3 \\
\bottomrule
\end{tabular}
\end{center}
工作 6 填了時段 $4$ 以後，\proc{Union}$(4, 3)$：$4$ 那組只有 $1$ 個，$3$ 那組有 $2$
個，所以 $4$ 接到 $3$ 底下，$\mathit{latest}[3] = 2$。 時段 $1$ 到 $5$ 依序是工作
$3, 1, 4, 6, 7$，總收費 $30 + 40 + 55 + 45 + 10 = 180$。 用範例~\ref{ex:gr-deadline}
跑這組輸入，印出的父節點陣列和上表相同。

\item
\begin{enumerate}
\item 金額 $1$ 到 $11$ 貪心都對（例如 $8 = 4 + 4$、$11 = 9 + 1 + 1$）。 金額 $12$：
貪心拿 $9 + 1 + 1 + 1$，$4$ 枚；$4 + 4 + 4$ 只要 $3$ 枚。 最小的反例是 $12$，也符合
Kozen--Zaks 的界 $12 < 4 + 9 = 13$。
\item 只要檢查 $1$ 到 $59$。 最佳湊法裡 $1$ 元最多 $4$ 枚、$5$ 元最多 $1$ 枚、$10$ 元最多
$4$ 枚（否則可以換成更大的一枚，枚數變少），這三種加起來最多 $49$ 元；所以金額
$\ge 50$ 時最佳湊法一定用到 $50$ 元，扣掉後剩下的金額 $\ge 10$ 時一定用到 $10$ 元，依此類
推，正是貪心的選擇。 範圍內貪心都最少，這組面額是標準的。
\end{enumerate}

\item 對 $i$ 歸納。 $i = 1$：$g_1$ 是所有申請裡結束最早的，$f(g_1) \le f(o_1)$。 假設
$f(g_{i-1}) \le f(o_{i-1})$。 $o_i$ 和 $o_{i-1}$ 相容而且在它之後，
$s(o_i) \ge f(o_{i-1}) \ge f(g_{i-1})$，所以選完 $g_{i-1}$ 時，$o_i$ 仍然可以選。 貪心
接下來選的 $g_i$，是開始時刻 $\ge f(g_{i-1})$ 的申請裡結束最早的，所以
$f(g_i) \le f(o_i)$。

假設 $m > k$。 由上面的結果 $f(g_k) \le f(o_k) \le s(o_{k+1})$，所以 $o_{k+1}$ 在
$g_k$ 之後仍然可以選，而它的結束時刻 $f(o_{k+1}) > s(o_{k+1}) \ge f(g_k)$，排在
$g_k$ 後面。 貪心掃到它的時候 $\mathit{last} = f(g_k) \le s(o_{k+1})$，會把它選進來，
和「貪心只選了 $k$ 筆」矛盾。 所以 $k \ge m$，貪心選的筆數不少於任何一組相容的申請。
這種寫法叫\textbf{逐步領先}（greedy stays ahead）：不改動最佳解，直接證明貪心每一步
都不落後。

\item (a)、(c)、(e)。 (a) 是圖的擬陣（第~\ref{sec:gr-matroid}~節）。 (b)：$A =
\{[0, 5)\}$、$B = \{[0, 1), [2, 3)\}$，$B$ 的區間都和 $[0, 5)$ 重疊，交換性不成立。
(c) 在第~\ref{sec:gr-matroid}~節說明過。 (d)：路徑 $a - b - c - d$，$A = \{bc\}$、
$B = \{ab, cd\}$，交換性不成立。 (e)：$|A| < |B| \le 3$ 時，$B \setminus A$ 的任何元素加進 $A$ 都不超過 $3$
個，交換性成立；這叫均勻擬陣。
\end{enumerate}
